\documentclass[12pt,a4paper]{article}
\usepackage{header} % Provided header format
\usepackage[utf8]{inputenc}
\usepackage[T5]{fontenc}
\usepackage[vietnamese,english]{babel}
\usetikzlibrary{calc,patterns}

% ---------------------------------------------------------------
% Helper macros
% ---------------------------------------------------------------
% Origin marker: \org{3} -> 3 with an up-arrow underneath
\newcommand{\org}[1]{\underset{\uparrow}{#1}}

% Reserved (dashed) box for parts that are written by students / other members.
% Usage: \writespace[height]{hint text}
\newcommand{\writespace}[2][4.5cm]{%
  \begin{tcolorbox}[enhanced, frame hidden, colback=gray!4,
      borderline={0.6pt}{0pt}{gray!70, dashed},
      height=#1, valign=center, halign=center,
      fontupper=\itshape\small\color{gray}]
    #2
  \end{tcolorbox}}

% Scilab code listing macro
\newcommand{\scilabfile}[2]{%
  \lstinputlisting[language=Scilab, caption={#2}]{#1}}

\begin{document}

\begin{center}
    \textbf{Digital Signal Processing \hfill HK261 \hfill CO2036}\\[2cm]

    \huge\textbf{DIGITAL SIGNAL PROCESSING LAB REPORT}\\[0.5cm]
    \Large\textbf{Lab 5 -- The Z-Transform, Inverse Z-Transform, and LTI System Analysis}\\[1cm]

    \large Âu Lê Quang Trí - Student ID: 2453300\\[0.5cm]
    \large Nguyễn Lê Khải Tú - Student ID: 2453384\\[0.5cm]
    \vspace{0.5cm}

    \large \today\\[2cm]
\end{center}

\section*{Overview}
This laboratory report addresses the analytical foundations and applications of the bilateral $\mathcal{Z}$-transform, the unilateral $\mathcal{Z}^+$-transform, and the inverse $\mathcal{Z}$-transform ($\mathcal{Z}^{-1}$) for discrete-time signal and linear time-invariant (LTI) system analysis. 

The report is structured in two major sections:
\begin{enumerate}
    \item \textbf{Part I: Core Laboratory Exercises (Lab Manual):} Focuses on time-domain convolution via algebraic $\mathcal{Z}$-domain multiplication, block diagram representations of recursive discrete-time difference equations, impulse and step responses, and fundamental operational properties (complex conjugation, real/imaginary sequence decomposition, and complex frequency modulation).
    \item \textbf{Part II: Additional Exercises (Reference Document -- Proakis \& Manolakis, Ch.~3):} Examines $\mathcal{Z}$-transform derivations with detailed Region of Convergence (ROC) specifications, accumulation and self-convolution properties, rational inverse $\mathcal{Z}$-transforms for causal systems (including pole-zero pattern reconstruction), and parametric pole-zero modeling of damped sinusoidal sequences.
\end{enumerate}

\tableofcontents

\newpage

% =========================================================================
% REFERENCE TABLES
% =========================================================================
\section*{Reference Tables}

\begin{table}[H]
\centering
\caption{Properties of the $z$-Transform}
\label{tab:z_properties}
{\small
\setlength{\tabcolsep}{5pt}
\renewcommand{\arraystretch}{1.35}
\begin{tabular}{>{\raggedright\arraybackslash}p{3.0cm}>{\raggedright\arraybackslash}p{3.4cm}>{\raggedright\arraybackslash}p{4.5cm}>{\raggedright\arraybackslash}p{4.3cm}}
\hline
\textbf{Property} & \textbf{Time Domain} & \textbf{$z$-Domain} & \textbf{ROC} \\
\hline
Notation & $x(n)$ & $X(z)$ & $\text{ROC}: r_2 < |z| < r_1$ \\
 & $x_1(n)$ & $X_1(z)$ & $\text{ROC}_1$ \\
 & $x_2(n)$ & $X_2(z)$ & $\text{ROC}_2$ \\
Linearity & $a_1 x_1(n) + a_2 x_2(n)$ & $a_1 X_1(z) + a_2 X_2(z)$ & At least the intersection of $\text{ROC}_1$ and $\text{ROC}_2$ \\
Time shifting & $x(n - k)$ & $z^{-k}X(z)$ & That of $X(z)$, except $z = 0$ if $k > 0$ and $z = \infty$ if $k < 0$ \\
Scaling in the $z$-domain & $a^n x(n)$ & $X(a^{-1}z)$ & $|a|r_2 < |z| < |a|r_1$ \\
Time reversal & $x(-n)$ & $X(z^{-1})$ & $\frac{1}{r_1} < |z| < \frac{1}{r_2}$ \\[3pt]
Conjugation & $x^*(n)$ & $X^*(z^*)$ & $\text{ROC}$ \\
Real part & $\text{Re}\{x(n)\}$ & $\frac{1}{2}\left[X(z) + X^*(z^*)\right]$ & Includes $\text{ROC}$ \\[2pt]
Imaginary part & $\text{Im}\{x(n)\}$ & $\frac{1}{2j}\left[X(z) - X^*(z^*)\right]$ & Includes $\text{ROC}$ \\[3pt]
Differentiation in the $z$-domain & $nx(n)$ & $-z\frac{dX(z)}{dz}$ & $r_2 < |z| < r_1$ \\[3pt]
Convolution & $x_1(n) * x_2(n)$ & $X_1(z)X_2(z)$ & At least, the intersection of $\text{ROC}_1$ and $\text{ROC}_2$ \\
Correlation & $r_{x_1 x_2}(l) = x_1(l) * x_2(-l)$ & $R_{x_1 x_2}(z) = X_1(z)X_2(z^{-1})$ & At least, the intersection of $\text{ROC}$ of $X_1(z)$ and $X_2(z^{-1})$ \\
Initial value theorem & If $x(n)$ causal & $x(0) = \lim_{z \to \infty} X(z)$ & \\[2pt]
Multiplication & $x_1(n)x_2(n)$ & $\frac{1}{2\pi j}\oint_C X_1(v)X_2\left(\frac{z}{v}\right)v^{-1}dv$ & At least, $r_{1l}r_{2l} < |z| < r_{1u}r_{2u}$ \\[4pt]
Parseval's relation & \multicolumn{2}{l}{$\displaystyle\sum_{n=-\infty}^{\infty} x_1(n)x_2^*(n) = \frac{1}{2\pi j}\oint_C X_1(v)X_2^*(1/v^*)v^{-1}dv$} & \\[6pt]
\hline
\end{tabular}
}
\end{table}

\newpage

\begin{table}[H]
\centering
\caption{Some Common $z$-Transform Pairs}
\label{tab:z_pairs}
\renewcommand{\arraystretch}{1.3}
\begin{tabular}{cccc}
\hline
\rule[-1.5ex]{0pt}{4.2ex}\textbf{\#} & \textbf{Signal, $x(n)$} & \textbf{$z$-Transform, $X(z)$} & \textbf{ROC} \\
\hline
\rule[-1.8ex]{0pt}{4.5ex}1 & $\delta(n)$ & $1$ & All $z$ \\
\rule[-3.2ex]{0pt}{7.0ex}2 & $u(n)$ & $\dfrac{1}{1 - z^{-1}}$ & $|z| > 1$ \\
\rule[-3.2ex]{0pt}{7.0ex}3 & $a^n u(n)$ & $\dfrac{1}{1 - az^{-1}}$ & $|z| > |a|$ \\
\rule[-3.2ex]{0pt}{7.0ex}4 & $na^n u(n)$ & $\dfrac{az^{-1}}{(1 - az^{-1})^2}$ & $|z| > |a|$ \\
\rule[-3.2ex]{0pt}{7.0ex}5 & $-a^n u(-n - 1)$ & $\dfrac{1}{1 - az^{-1}}$ & $|z| < |a|$ \\
\rule[-3.2ex]{0pt}{7.0ex}6 & $-na^n u(-n - 1)$ & $\dfrac{az^{-1}}{(1 - az^{-1})^2}$ & $|z| < |a|$ \\
\rule[-3.2ex]{0pt}{7.0ex}7 & $(\cos\omega_0 n)u(n)$ & $\dfrac{1 - z^{-1}\cos\omega_0}{1 - 2z^{-1}\cos\omega_0 + z^{-2}}$ & $|z| > 1$ \\
\rule[-3.2ex]{0pt}{7.0ex}8 & $(\sin\omega_0 n)u(n)$ & $\dfrac{z^{-1}\sin\omega_0}{1 - 2z^{-1}\cos\omega_0 + z^{-2}}$ & $|z| > 1$ \\
\rule[-3.2ex]{0pt}{7.0ex}9 & $(a^n \cos\omega_0 n)u(n)$ & $\dfrac{1 - az^{-1}\cos\omega_0}{1 - 2az^{-1}\cos\omega_0 + a^2 z^{-2}}$ & $|z| > |a|$ \\
\rule[-3.2ex]{0pt}{7.0ex}10 & $(a^n \sin\omega_0 n)u(n)$ & $\dfrac{az^{-1}\sin\omega_0}{1 - 2az^{-1}\cos\omega_0 + a^2 z^{-2}}$ & $|z| > |a|$ \\
\hline
\end{tabular}
\end{table}

\newpage

% =========================================================================
% PART I: CORE LABORATORY EXERCISES
% =========================================================================
\section{Part I: Core Laboratory Exercises}

% -------------------------------------------------------------------------
% Exercise 1
% -------------------------------------------------------------------------
\subsection{Exercise 1: Impulse Response via Z-Transform}
\textbf{Problem Statement:} Use the $\mathcal{Z}$-transform to find the impulse response $h(n)$ of the system represented by the following input-output difference equation:
\begin{equation}
    y(n) - y(n-2) = x(n).
    \label{eq:ex1_diff}
\end{equation}

\textbf{Solution:}
Take the $z$-transform of both sides of the difference equation using the linearity and time-shifting properties from Table~\ref{tab:z_properties}:
\[
Y(z) - z^{-2}Y(z) = X(z)
\]
Factor out $Y(z)$ on the left side:
\[
Y(z)(1 - z^{-2}) = X(z)
\]
The system function $H(z)$ is the ratio of $Y(z)$ to $X(z)$:
\[
H(z) = \frac{Y(z)}{X(z)} = \frac{1}{1 - z^{-2}}
\]
To invert $H(z)$ by partial-fraction expansion, multiply the numerator and denominator by $z^2$ to eliminate negative powers of $z$:
\[
H(z) = \frac{z^2}{z^2 - 1} = \frac{z^2}{(z - 1)(z + 1)}
\]
Divide both sides by $z$:
\[
\frac{H(z)}{z} = \frac{z}{(z - 1)(z + 1)} = \frac{A_1}{z - 1} + \frac{A_2}{z + 1}
\]
Find the coefficients $A_1$ and $A_2$:
\[
A_1 = \left. \frac{(z - 1)H(z)}{z} \right|_{z = 1} = \left. \frac{z}{z + 1} \right|_{z = 1} = \frac{1}{2}
\]
\[
A_2 = \left. \frac{(z + 1)H(z)}{z} \right|_{z = -1} = \left. \frac{z}{z - 1} \right|_{z = -1} = \frac{-1}{-2} = \frac{1}{2}
\]
Multiply both sides by $z$ and rewrite the expression in powers of $z^{-1}$:
\[
H(z) = \frac{1}{2}\frac{1}{1 - z^{-1}} + \frac{1}{2}\frac{1}{1 - (-1)z^{-1}}
\]
The poles of $H(z)$ are $p_1 = 1$ and $p_2 = -1$, which both have magnitude $1$. For a causal system, the region of convergence (ROC) is $|z| > 1$, and Table~\ref{tab:z_pairs} gives the causal impulse response:
\[
h(n) = \frac{1}{2}u(n) + \frac{1}{2}(-1)^n u(n) = \frac{1}{2}\left[1 + (-1)^n\right]u(n)
\]
If the system is anticausal with ROC $|z| < 1$, the impulse response is:
\[
h(n) = -\frac{1}{2}\left[1 + (-1)^n\right]u(-n - 1)
\]

% -------------------------------------------------------------------------
% Exercise 2
% -------------------------------------------------------------------------
\subsection{Exercise 2: Convolution via Z and Inverse Z Transforms}
\textbf{Problem Statement:} Use the forward $\mathcal{Z}$-transform and inverse $\mathcal{Z}$-transform ($\mathcal{Z}^{-1}$) to compute the convolution $y(n) = x_1(n) * x_2(n)$ for each of the following signal pairs:
\begin{enumerate}[label=(\alph*)]
    \item $x_1(n) = \{\org{1}, 2, 3, 4, 5\}$ \quad and \quad $x_2(n) = \{\org{1}, 1, 1\}$
    \item $x_1(n) = \left(\dfrac{1}{5}\right)^n u(n)$ \quad and \quad $x_2(n) = 2^n u(n)$
    \item $x_1(n) = n u(n)$ \quad and \quad $x_2(n) = 2^n u(n-1)$
\end{enumerate}

\subsubsection*{a)}
The input signals are finite-duration causal sequences:
\begin{align*}
    x_1(n) &= \delta(n) + 2\delta(n-1) + 3\delta(n-2) + 4\delta(n-3) + 5\delta(n-4), \\
    x_2(n) &= \delta(n) + \delta(n-1) + \delta(n-2).
\end{align*}
Applying the forward $\mathcal{Z}$-transform by direct definition:
\begin{align*}
    X_1(z) &= \sum_{n=0}^{4} x_1(n) z^{-n} = 1 + 2z^{-1} + 3z^{-2} + 4z^{-3} + 5z^{-4}, \quad \text{ROC: } |z| > 0, \\
    X_2(z) &= \sum_{n=0}^{2} x_2(n) z^{-n} = 1 + z^{-1} + z^{-2}, \quad \text{ROC: } |z| > 0.
\end{align*}
By the convolution theorem in the $\mathcal{Z}$-domain, $Y(z) = X_1(z) X_2(z)$:
\begin{align*}
    Y(z) &= (1 + 2z^{-1} + 3z^{-2} + 4z^{-3} + 5z^{-4})(1 + z^{-1} + z^{-2}) \\
         &= 1(1 + z^{-1} + z^{-2}) + 2z^{-1}(1 + z^{-1} + z^{-2}) + 3z^{-2}(1 + z^{-1} + z^{-2}) \\
         &\quad + 4z^{-3}(1 + z^{-1} + z^{-2}) + 5z^{-4}(1 + z^{-1} + z^{-2}).
\end{align*}
Collecting terms by powers of $z^{-1}$:
\begin{align*}
    z^0:    &\quad 1 = 1, \\
    z^{-1}: &\quad 1 + 2 = 3, \\
    z^{-2}: &\quad 1 + 2 + 3 = 6, \\
    z^{-3}: &\quad 2 + 3 + 4 = 9, \\
    z^{-4}: &\quad 3 + 4 + 5 = 12, \\
    z^{-5}: &\quad 4 + 5 = 9, \\
    z^{-6}: &\quad 5 = 5.
\end{align*}
Thus, the product transform is
\begin{equation}
    Y(z) = 1 + 3z^{-1} + 6z^{-2} + 9z^{-3} + 12z^{-4} + 9z^{-5} + 5z^{-6}, \quad \text{ROC: } |z| > 0.
\end{equation}
Taking the inverse $\mathcal{Z}$-transform by inspecting polynomial coefficients:
\begin{equation}
    y(n) = \{\org{1}, 3, 6, 9, 12, 9, 5\}.
\end{equation}

\subsubsection*{b)}
Given the infinite-duration right-sided exponential sequences:
\begin{align*}
    x_1(n) &= \left(\frac{1}{5}\right)^n u(n) \implies X_1(z) = \frac{1}{1 - \frac{1}{5}z^{-1}}, \quad \text{ROC}_1: |z| > \frac{1}{5}, \\
    x_2(n) &= 2^n u(n) \implies X_2(z) = \frac{1}{1 - 2z^{-1}}, \quad \text{ROC}_2: |z| > 2.
\end{align*}
The $\mathcal{Z}$-transform of the convolution is
\begin{equation}
    Y(z) = X_1(z) X_2(z) = \frac{1}{\left(1 - \frac{1}{5}z^{-1}\right)(1 - 2z^{-1})}.
\end{equation}
The overall ROC is the intersection $\text{ROC}_1 \cap \text{ROC}_2 = \{z \in \mathbb{C} : |z| > 2\}$.

We perform a partial fraction expansion in terms of $z^{-1}$:
\begin{equation}
    Y(z) = \frac{A}{1 - \frac{1}{5}z^{-1}} + \frac{B}{1 - 2z^{-1}}.
\end{equation}
Using Heaviside's residue method:
\begin{align*}
    A &= \left. \frac{1}{1 - 2z^{-1}} \right|_{z^{-1} = 5} = \frac{1}{1 - 10} = -\frac{1}{9}, \\
    B &= \left. \frac{1}{1 - \frac{1}{5}z^{-1}} \right|_{z^{-1} = 1/2} = \frac{1}{1 - \frac{1}{10}} = \frac{10}{9}.
\end{align*}
Hence:
\begin{equation}
    Y(z) = -\frac{1}{9} \frac{1}{1 - \frac{1}{5}z^{-1}} + \frac{10}{9} \frac{1}{1 - 2z^{-1}}, \quad |z| > 2.
\end{equation}
Because the ROC $|z| > 2$ is exterior to both poles at $z = 1/5$ and $z = 2$, both terms correspond to causal right-sided sequences:
\begin{equation}
    y(n) = \left[ -\frac{1}{9}\left(\frac{1}{5}\right)^n + \frac{10}{9} (2^n) \right] u(n) = \frac{1}{9}\left[ 10 \cdot 2^n - \left(\frac{1}{5}\right)^n \right] u(n).
\end{equation}

\subsubsection*{c)}
The input sequences are:
\begin{equation*}
    x_1(n) = n u(n), \qquad x_2(n) = 2^n u(n-1).
\end{equation*}
Recall the differentiation property in the $\mathcal{Z}$-domain: $\mathcal{Z}\{n v(n)\} = -z \frac{d}{dz} V(z)$. Since $\mathcal{Z}\{u(n)\} = \frac{1}{1 - z^{-1}}$ for $|z| > 1$:
\begin{equation}
    X_1(z) = -z \frac{d}{dz}\left(\frac{1}{1 - z^{-1}}\right) = -z \left( -\frac{z^{-2}}{(1 - z^{-1})^2} \right) = \frac{z^{-1}}{(1 - z^{-1})^2}, \quad |z| > 1.
\end{equation}
For $x_2(n) = 2^n u(n-1) = 2 \cdot 2^{n-1} u(n-1)$:
\begin{equation}
    X_2(z) = \sum_{n=1}^{\infty} 2^n z^{-n} = 2z^{-1} \sum_{m=0}^{\infty} (2z^{-1})^m = \frac{2z^{-1}}{1 - 2z^{-1}}, \quad |z| > 2.
\end{equation}
The product transform is
\begin{equation}
    Y(z) = X_1(z) X_2(z) = \frac{2z^{-2}}{(1 - z^{-1})^2(1 - 2z^{-1})}, \quad |z| > 2.
\end{equation}
Letting $w = z^{-1}$, we expand $\frac{2w^2}{(1 - w)^2(1 - 2w)}$:
\begin{equation}
    \frac{2w^2}{(1 - w)^2(1 - 2w)} = \frac{A_1}{1 - w} + \frac{A_2}{(1 - w)^2} + \frac{B}{1 - 2w}.
\end{equation}
Evaluating the coefficients:
\begin{align*}
    B   &= \left. \frac{2w^2}{(1 - w)^2} \right|_{w = 1/2} = \frac{2(1/4)}{(1/2)^2} = 2, \\
    A_2 &= \left. \frac{2w^2}{1 - 2w} \right|_{w = 1} = \frac{2(1)}{1 - 2} = -2.
\end{align*}
Setting $w = 0$ on both sides:
\begin{equation*}
    0 = A_1 + A_2 + B \implies A_1 = -A_2 - B = -(-2) - 2 = 0.
\end{equation*}
Therefore:
\begin{equation}
    Y(z) = \frac{2}{1 - 2z^{-1}} - \frac{2}{(1 - z^{-1})^2}, \quad |z| > 2.
\end{equation}
Using the standard transform pairs:
\begin{align*}
    \mathcal{Z}^{-1}\left\{ \frac{1}{1 - 2z^{-1}} \right\} &= 2^n u(n), \\
    \mathcal{Z}^{-1}\left\{ \frac{1}{(1 - z^{-1})^2} \right\} &= (n+1)u(n).
\end{align*}
Taking the inverse $\mathcal{Z}$-transform yields:
\begin{align}
    y(n) &= \left[ 2 \cdot 2^n - 2(n+1) \right] u(n) \nonumber \\
         &= \left( 2^{n+1} - 2n - 2 \right) u(n).
\end{align}
\textit{Verification:} For $n=0$: $y(0) = 2^1 - 0 - 2 = 0 = x_1(0)x_2(0)$. For $n=1$: $y(1) = 2^2 - 2 - 2 = 0$. For $n=2$: $y(2) = 2^3 - 4 - 2 = 2 = x_1(1)x_2(1) = 1 \cdot 2 = 2$. The analytical result agrees perfectly.

% -------------------------------------------------------------------------
% Exercise 3
% -------------------------------------------------------------------------
\subsection{Exercise 3: Inverse Z-Transform for All Possible Sequences}
\textbf{Problem Statement:} Find all possible discrete-time sequences $x(n)$ corresponding to the following $\mathcal{Z}$-transforms based on all admissible Regions of Convergence:
\begin{align*}
    X_1(z) &= \frac{1}{2 - 3z^{-1} + z^{-2}}, \\
    X_2(z) &= \frac{1 + 2z^{-1} + z^{-2}}{1 + 4z^{-1} + 4z^{-2}}, \\
    X_3(z) &= \frac{1 + z^{-1}}{(1 - 0.5z^{-1})^2 (1 - 0.3z^{-1})}, \quad \text{with specified ROC: } 0.3 < |z| < 0.5.
\end{align*}

\subsubsection*{a)}
Multiply the numerator and denominator by $z^2$ to express $X_1(z)$ in positive powers of $z$:
\[
X_1(z) = \frac{z^2}{2z^2 - 3z + 1}
\]
Divide both sides by $z$ and factor the denominator:
\[
\begin{aligned}
\frac{X_1(z)}{z} &= \frac{z}{(2z - 1)(z - 1)} = \frac{z}{2(z - 0.5)(z - 1)} \\
&= \frac{A_1}{z - 0.5} + \frac{A_2}{z - 1}
\end{aligned}
\]
The poles are $p_1 = 0.5$ and $p_2 = 1$. Calculate the coefficients $A_1$ and $A_2$:
\[
A_1 = \left. \frac{(z - 0.5)X_1(z)}{z} \right|_{z = 0.5} = \left. \frac{z}{2(z - 1)} \right|_{z = 0.5} = \frac{0.5}{2(0.5 - 1)} = -\frac{1}{2}
\]
\[
A_2 = \left. \frac{(z - 1)X_1(z)}{z} \right|_{z = 1} = \left. \frac{z}{2(z - 0.5)} \right|_{z = 1} = \frac{1}{2(1 - 0.5)} = 1
\]
Multiply by $z$ and write in powers of $z^{-1}$:
\[
X_1(z) = -\frac{1}{2}\frac{1}{1 - 0.5z^{-1}} + \frac{1}{1 - z^{-1}}
\]
Because the poles have magnitudes $|p_1| = 0.5$ and $|p_2| = 1$, there are three possible regions of convergence and three corresponding signals $x_1(n)$ from Table~\ref{tab:z_pairs}:
\begin{enumerate}
\item When the ROC is $|z| > 1$, the signal is causal:
\[
x_1(n) = \left[1 - \frac{1}{2}(0.5)^n\right]u(n)
\]
\item When the ROC is $|z| < 0.5$, the signal is anticausal:
\[
x_1(n) = \left[\frac{1}{2}(0.5)^n - 1\right]u(-n - 1)
\]
\item When the ROC is $0.5 < |z| < 1$, the signal is two-sided (the pole at $0.5$ gives a causal term and the pole at $1$ gives an anticausal term):
\[
x_1(n) = -\frac{1}{2}(0.5)^n u(n) - u(-n - 1)
\]
\end{enumerate}

\subsubsection*{b)}
Because the numerator and denominator have the same degree ($M = N = 2$), $X_2(z)$ is an improper rational function. Perform polynomial long division in reverse order to eliminate the $z^{-2}$ term in the numerator:
\[
1 + 2z^{-1} + z^{-2} = \frac{1}{4}(1 + 4z^{-1} + 4z^{-2}) + \left(\frac{3}{4} + z^{-1}\right)
\]
This separates $X_2(z)$ into a constant and a proper rational function $X_{\text{pr}}(z)$:
\[
X_2(z) = \frac{1}{4} + \frac{\frac{3}{4} + z^{-1}}{1 + 4z^{-1} + 4z^{-2}} = \frac{1}{4} + X_{\text{pr}}(z)
\]
Multiply the numerator and denominator of $X_{\text{pr}}(z)$ by $z^2$ and divide by $z$:
\[
\frac{X_{\text{pr}}(z)}{z} = \frac{\frac{3}{4}z + 1}{z^2 + 4z + 4} = \frac{\frac{3}{4}z + 1}{(z + 2)^2}
\]
This function has a double pole at $p = -2$. Write its partial-fraction expansion as:
\[
\frac{X_{\text{pr}}(z)}{z} = \frac{A_1}{z + 2} + \frac{A_2}{(z + 2)^2}
\]
Multiply both sides by $(z + 2)^2$:
\[
\frac{(z + 2)^2 X_{\text{pr}}(z)}{z} = \frac{3}{4}z + 1 = (z + 2)A_1 + A_2
\]
Evaluate at $z = -2$ to find $A_2$:
\[
A_2 = \left. \left(\frac{3}{4}z + 1\right) \right|_{z = -2} = -\frac{3}{2} + 1 = -\frac{1}{2}
\]
Differentiate with respect to $z$ and evaluate at $z = -2$ to find $A_1$:
\[
A_1 = \left. \frac{d}{dz}\left(\frac{3}{4}z + 1\right) \right|_{z = -2} = \frac{3}{4}
\]
Multiply by $z$ and rewrite $X_2(z)$ in terms of $z^{-1}$ to match the standard $z$-transform pairs in Table~\ref{tab:z_pairs}:
\[
\begin{aligned}
X_2(z) &= \frac{1}{4} + \frac{3}{4}\frac{1}{1 + 2z^{-1}} - \frac{1}{2}\frac{z^{-1}}{(1 + 2z^{-1})^2} \\
&= \frac{1}{4} + \frac{3}{4}\frac{1}{1 - (-2)z^{-1}} + \frac{1}{4}\frac{(-2)z^{-1}}{\left[1 - (-2)z^{-1}\right]^2}
\end{aligned}
\]
Since the only pole is at $p = -2$ with magnitude $|p| = 2$, there are two possible regions of convergence:
\begin{enumerate}
\item When the ROC is $|z| > 2$, the signal is causal:
\[
\begin{aligned}
x_2(n) &= \frac{1}{4}\delta(n) + \frac{3}{4}(-2)^n u(n) + \frac{1}{4}n(-2)^n u(n) \\
&= \frac{1}{4}\delta(n) + \frac{n + 3}{4}(-2)^n u(n)
\end{aligned}
\]
\item When the ROC is $|z| < 2$, the signal is anticausal for $n \neq 0$:
\[
\begin{aligned}
x_2(n) &= \frac{1}{4}\delta(n) - \frac{3}{4}(-2)^n u(-n - 1) - \frac{1}{4}n(-2)^n u(-n - 1) \\
&= \frac{1}{4}\delta(n) - \frac{n + 3}{4}(-2)^n u(-n - 1)
\end{aligned}
\]
\end{enumerate}

\subsubsection*{c)}
Multiply the numerator and denominator by $z^3$ to remove negative powers of $z$:
\[
X_3(z) = \frac{z^2(z + 1)}{(z - 0.5)^2(z - 0.3)}
\]
Divide both sides by $z$:
\[
\frac{X_3(z)}{z} = \frac{z(z + 1)}{(z - 0.3)(z - 0.5)^2} = \frac{A_1}{z - 0.3} + \frac{A_2}{z - 0.5} + \frac{A_3}{(z - 0.5)^2}
\]
The transform has a simple pole at $p_1 = 0.3$ and a double pole at $p_2 = 0.5$. Find $A_1$ by multiplying by $(z - 0.3)$ and evaluating at $z = 0.3$:
\[
A_1 = \left. \frac{z(z + 1)}{(z - 0.5)^2} \right|_{z = 0.3} = \frac{0.3(1.3)}{(-0.2)^2} = \frac{0.39}{0.04} = \frac{39}{4} = 9.75
\]
Find $A_3$ by multiplying by $(z - 0.5)^2$ and evaluating at $z = 0.5$:
\[
A_3 = \left. \frac{z(z + 1)}{z - 0.3} \right|_{z = 0.5} = \frac{0.5(1.5)}{0.2} = \frac{0.75}{0.2} = \frac{15}{4} = 3.75
\]
To find $A_2$, differentiate the product
\[
\frac{(z - 0.5)^2 X_3(z)}{z} = \frac{z^2 + z}{z - 0.3}
\]
with respect to $z$ and evaluate at $z = 0.5$:
\[
\begin{aligned}
A_2 &= \left. \frac{d}{dz}\left(\frac{z^2 + z}{z - 0.3}\right) \right|_{z = 0.5} = \left. \frac{(2z + 1)(z - 0.3) - (z^2 + z)}{(z - 0.3)^2} \right|_{z = 0.5} \\
&= \left. \frac{z^2 - 0.6z - 0.3}{(z - 0.3)^2} \right|_{z = 0.5} = \frac{0.25 - 0.30 - 0.30}{(0.2)^2} = -\frac{35}{4} = -8.75
\end{aligned}
\]
Multiply by $z$ and use the identity
\[
\frac{z^{-1}}{(1 - 0.5z^{-1})^2} = 2\frac{0.5z^{-1}}{(1 - 0.5z^{-1})^2}
\]
to write $X_3(z)$ in standard $z^{-1}$ terms:
\[
X_3(z) = \frac{39}{4}\frac{1}{1 - 0.3z^{-1}} - \frac{35}{4}\frac{1}{1 - 0.5z^{-1}} + \frac{15}{2}\frac{0.5z^{-1}}{(1 - 0.5z^{-1})^2}
\]
The given ROC is $0.3 < |z| < 0.5$, which is the intersection of $|z| > 0.3$ and $|z| < 0.5$. From Table~\ref{tab:z_pairs}, the pole at $0.3$ gives a causal component, and the double pole at $0.5$ gives anticausal components:
\[
\begin{aligned}
x_3(n) &= \frac{39}{4}(0.3)^n u(n) + \frac{35}{4}(0.5)^n u(-n - 1) - \frac{15}{2}n(0.5)^n u(-n - 1) \\
&= 9.75(0.3)^n u(n) + (8.75 - 7.5n)(0.5)^n u(-n - 1)
\end{aligned}
\]

% -------------------------------------------------------------------------
% Exercise 4
% -------------------------------------------------------------------------
\subsection{Exercise 4: First-Order LTI System Analysis and Implementation}
\textbf{Problem Statement:} An LTI system is described by the difference equation:
\begin{equation}
    y(n) = 0.7 y(n-1) + x(n).
    \label{eq:ex4_diff}
\end{equation}
\begin{enumerate}[label=(\alph*)]
    \item Draw the block diagram representation of the system.
    \item Determine the unit impulse response $h(n)$.
    \item Determine the total output $y(n)$ when the input is the unit step sequence $x(n) = u(n)$.
\end{enumerate}

\subsubsection{Block Diagram Representation}
The difference equation specifies that the current output sample $y(n)$ is formed by summing the current input sample $x(n)$ with the previous output sample $y(n-1)$ scaled by the feedback coefficient $a_1 = 0.7$. The block diagram is illustrated in Figure~\ref{fig:ex4_block_diagram}.

\begin{figure}[H]
    \centering
    \begin{tikzpicture}[auto, >=latex', node distance=2.2cm, thick]
        % Nodes
        \node [coordinate] (input) {};
        \node [circle, draw, minimum size=0.8cm, inner sep=0pt, right of=input] (sum) {$\mathbf{+}$};
        \node [coordinate, right of=sum, node distance=2.2cm] (branch) {};
        \node [coordinate, right of=branch, node distance=1.8cm] (output) {};
        
        \node [draw, rectangle, minimum width=1.4cm, minimum height=1.0cm, below of=branch, node distance=1.8cm] (delay) {$z^{-1}$};
        \node [draw, rectangle, minimum width=1.4cm, minimum height=1.0cm, left of=delay, node distance=2.2cm] (gain) {$0.7$};

        % Connections
        \draw [->] (input) -- node [above] {$x(n)$} (sum);
        \draw [-]  (sum) -- (branch);
        \draw [->] (branch) -- node [above] {$y(n)$} (output);
        \draw [->] (branch) -- (delay);
        \draw [->] (delay) -- node [above] {$y(n-1)$} (gain);
        \draw [->] (gain.north) -- node [left] {$+$} (sum.south);
    \end{tikzpicture}
    \caption{Direct-form recursive block diagram for the first-order system $y(n) = 0.7y(n-1) + x(n)$.}
    \label{fig:ex4_block_diagram}
\end{figure}

\subsubsection{Determination of Impulse Response h(n)}
Applying the bilateral $\mathcal{Z}$-transform to Eq.~\eqref{eq:ex4_diff} with initial rest conditions:
\begin{equation}
    Y(z) = 0.7 z^{-1} Y(z) + X(z) \implies Y(z)(1 - 0.7 z^{-1}) = X(z).
\end{equation}
The transfer function $H(z)$ is given by
\begin{equation}
    H(z) = \frac{Y(z)}{X(z)} = \frac{1}{1 - 0.7 z^{-1}}.
\end{equation}
The pole is located at $z = 0.7$. For a causal and stable discrete-time system, the Region of Convergence is exterior to the pole:
\begin{equation}
    \text{ROC}_H: |z| > 0.7.
\end{equation}
Taking the inverse $\mathcal{Z}$-transform:
\begin{equation}
    h(n) = (0.7)^n u(n).
\end{equation}

\subsubsection{Determination of Step Response y(n) for x(n) = u(n)}
When $x(n) = u(n)$, its $\mathcal{Z}$-transform is
\begin{equation}
    X(z) = \frac{1}{1 - z^{-1}}, \quad \text{ROC}_X: |z| > 1.
\end{equation}
The system output in the $\mathcal{Z}$-domain is
\begin{equation}
    Y(z) = H(z) X(z) = \frac{1}{(1 - 0.7 z^{-1})(1 - z^{-1})}, \quad \text{ROC}_Y: |z| > 1.
\end{equation}
We expand $Y(z)$ using partial fractions:
\begin{equation}
    Y(z) = \frac{A}{1 - 0.7 z^{-1}} + \frac{B}{1 - z^{-1}}.
\end{equation}
The residues are computed as follows:
\begin{align*}
    A &= \left. \frac{1}{1 - z^{-1}} \right|_{z^{-1} = 1/0.7} = \frac{1}{1 - \frac{10}{7}} = -\frac{7}{3}, \\
    B &= \left. \frac{1}{1 - 0.7 z^{-1}} \right|_{z^{-1} = 1} = \frac{1}{1 - 0.7} = \frac{1}{0.3} = \frac{10}{3}.
\end{align*}
Thus:
\begin{equation}
    Y(z) = -\frac{7}{3} \frac{1}{1 - 0.7 z^{-1}} + \frac{10}{3} \frac{1}{1 - z^{-1}}, \quad |z| > 1.
\end{equation}
Inverting back to the time domain:
\begin{align}
    y(n) &= \left[ -\frac{7}{3}(0.7)^n + \frac{10}{3} \right] u(n) \nonumber \\
         &= \frac{10}{3}\left[ 1 - (0.7)(0.7)^n \right] u(n) = \frac{10}{3}\left[ 1 - (0.7)^{n+1} \right] u(n).
\end{align}
\textit{Numerical verification:}
\begin{itemize}
    \item At $n=0$: $y(0) = \frac{10}{3}(1 - 0.7) = 1$. From difference equation: $y(0) = 0.7(0) + 1 = 1$.
    \item At $n=1$: $y(1) = \frac{10}{3}(1 - 0.49) = \frac{10}{3}(0.51) = 1.7$. From difference equation: $y(1) = 0.7(1) + 1 = 1.7$.
    \item Steady-state as $n \to \infty$: $y(\infty) = \frac{10}{3} \approx 3.333$, in exact accordance with the DC gain $H(1) = \frac{1}{1 - 0.7} = \frac{10}{3}$.
\end{itemize}

% -------------------------------------------------------------------------
% Exercise 5
% -------------------------------------------------------------------------
\subsection{Exercise 5: Second-Order LTI System with Initial Conditions}
\textbf{Problem Statement:} Given the LTI system represented by the difference equation:
\begin{equation}
    y(n) = 2y(n-1) - 3y(n-2) + x(n) + x(n-1).
    \label{eq:ex5_diff}
\end{equation}
\begin{enumerate}[label=(\alph*)]
    \item Draw the block diagram representation of the system.
    \item Determine the unit impulse response $h(n)$.
    \item Determine the zero-input response $y_{zi}(n)$ when initial conditions are $y(-1) = 1$ and $y(-2) = 1$.
    \item Determine the zero-state response $y_{zs}(n)$ when the input is $x(n) = 2^n u(n)$.
\end{enumerate}

\subsubsection*{a)}
The difference equation specifies that the current output sample $y(n)$ is formed by summing the current input $x(n)$, the delayed input $x(n-1)$, the delayed output $y(n-1)$ scaled by $2$, and the twice-delayed output $y(n-2)$ scaled by $-3$. The direct-form block diagram is shown in Figure~\ref{fig:ex5_block_diagram}.

\begin{figure}[H]
    \centering
    \begin{tikzpicture}[>=latex', thick]
        % Input, summing, and output nodes
        \coordinate (input) at (0, 0);
        \coordinate (bx) at (1.5, 0);
        \node[circle, draw, minimum size=0.8cm, inner sep=0pt] (sum1) at (5.6, 0) {$\mathbf{+}$};
        \node[circle, draw, minimum size=0.8cm, inner sep=0pt] (sum2) at (8.1, 0) {$\mathbf{+}$};
        \coordinate (by) at (12.6, 0);
        \coordinate (output) at (14.2, 0);

        % Feedforward delay block
        \node[draw, rectangle, minimum width=1.2cm, minimum height=0.9cm] (dx) at (3.1, -2.5) {$z^{-1}$};

        % Feedback summing, gain, and delay blocks
        \node[circle, draw, minimum size=0.8cm, inner sep=0pt] (sumfb) at (8.1, -2.5) {$\mathbf{+}$};
        \node[draw, rectangle, minimum width=1.2cm, minimum height=0.9cm] (g1) at (10.2, -2.5) {$2$};
        \node[draw, rectangle, minimum width=1.2cm, minimum height=0.9cm] (dy1) at (12.6, -1.25) {$z^{-1}$};
        \coordinate (by1) at (12.6, -2.5);
        \node[draw, rectangle, minimum width=1.2cm, minimum height=0.9cm] (dy2) at (12.6, -3.7) {$z^{-1}$};
        \node[draw, rectangle, minimum width=1.2cm, minimum height=0.9cm] (g2) at (10.2, -4.9) {$-3$};

        % Branch dots
        \fill (bx) circle (2pt);
        \fill (by) circle (2pt);
        \fill (by1) circle (2pt);

        % Forward signal path
        \draw[->] (input) -- node[above] {$x(n)$} (bx) -- (sum1.west);
        \draw[->] (sum1.east) -- (sum2.west);
        \draw[->] (sum2.east) -- (by) -- node[above] {$y(n)$} (output);

        % Feedforward delay path
        \draw[->] (bx) -- (1.5, -2.5) -- (dx.west);
        \draw[->] (dx.east) -- node[above] {$x(n-1)$} (5.6, -2.5) -- (sum1.south);

        % Feedback delay chain and gain paths
        \draw[->] (by) -- (dy1.north);
        \draw (dy1.south) -- (by1);
        \draw[->] (by1) -- node[above] {$y(n-1)$} (g1.east);
        \draw[->] (g1.west) -- (sumfb.east);
        \draw[->] (sumfb.north) -- (sum2.south);

        \draw[->] (by1) -- (dy2.north);
        \draw (dy2.south) -- (12.6, -4.9);
        \draw[->] (12.6, -4.9) -- node[above] {$y(n-2)$} (g2.east);
        \draw[->] (g2.west) -- (8.1, -4.9) -- (sumfb.south);
    \end{tikzpicture}
    \caption{Direct-form block diagram for the second-order system $y(n) = 2y(n-1) - 3y(n-2) + x(n) + x(n-1)$.}
    \label{fig:ex5_block_diagram}
\end{figure}

\subsubsection*{b)}
Rewrite the difference equation with all output terms on the left side:
\[
y(n) - 2y(n-1) + 3y(n-2) = x(n) + x(n-1)
\]
Take the $z$-transform of both sides for a relaxed causal system:
\[
Y(z)(1 - 2z^{-1} + 3z^{-2}) = X(z)(1 + z^{-1})
\]
Solve for the system function $H(z)$:
\[
H(z) = \frac{Y(z)}{X(z)} = \frac{1 + z^{-1}}{1 - 2z^{-1} + 3z^{-2}}
\]
The denominator has complex-conjugate poles at $p_{1,2} = 1 \pm j\sqrt{2}$. Compare the denominator $1 - 2z^{-1} + 3z^{-2}$ with the standard second-order form $1 - 2az^{-1}\cos\omega_0 + a^2 z^{-2}$ from Table~\ref{tab:z_pairs}:
\[
a^2 = 3 \implies a = \sqrt{3}
\]
\[
\begin{aligned}
2a\cos\omega_0 = 2 &\implies \cos\omega_0 = \frac{1}{\sqrt{3}}, \quad \sin\omega_0 = \sqrt{\frac{2}{3}} \\
&\implies \omega_0 = \arccos\left(\frac{1}{\sqrt{3}}\right) \approx 54.74^\circ
\end{aligned}
\]
This gives $a\cos\omega_0 = 1$ and $a\sin\omega_0 = \sqrt{2}$. Split the numerator $1 + z^{-1} = (1 - z^{-1}) + 2z^{-1}$ to match the cosine and sine transform pairs in Table~\ref{tab:z_pairs}:
\[
\begin{aligned}
H(z) &= \frac{1 - az^{-1}\cos\omega_0}{1 - 2az^{-1}\cos\omega_0 + a^2 z^{-2}} \\
&\quad + \sqrt{2}\frac{az^{-1}\sin\omega_0}{1 - 2az^{-1}\cos\omega_0 + a^2 z^{-2}}
\end{aligned}
\]
Invert both terms using Table~\ref{tab:z_pairs} to obtain the causal impulse response:
\[
h(n) = (\sqrt{3})^n \left[\cos(\omega_0 n) + \sqrt{2}\sin(\omega_0 n)\right]u(n)
\]
Equivalently, expand $\frac{H(z)}{z}$ in partial fractions at $p_1 = 1 + j\sqrt{2} = \sqrt{3}e^{j\omega_0}$:
\[
A_1 = \left. \frac{z + 1}{z - p_2} \right|_{z = 1 + j\sqrt{2}} = \frac{2 + j\sqrt{2}}{j2\sqrt{2}} = \frac{1}{2} - j\frac{\sqrt{2}}{2} = \frac{\sqrt{3}}{2}e^{-j\omega_0}
\]
This writes $h(n)$ as a single cosine sequence:
\[
h(n) = (\sqrt{3})^{n+1}\cos\left[\omega_0(n - 1)\right]u(n), \quad \omega_0 \approx 54.74^\circ
\]

\subsubsection*{c)}
Set the input $x(n) = 0$ in the difference equation:
\[
y_{zi}(n) - 2y_{zi}(n-1) + 3y_{zi}(n-2) = 0
\]
Take the one-sided $z$-transform of both sides using the time-delay shifting property:
\[
\begin{aligned}
Y_{zi}^+(z) &- 2\left[z^{-1}Y_{zi}^+(z) + y(-1)\right] \\
&+ 3\left[z^{-2}Y_{zi}^+(z) + y(-1)z^{-1} + y(-2)\right] = 0
\end{aligned}
\]
Collect the $Y_{zi}^+(z)$ terms on the left side and move the initial condition terms to the right side:
\[
Y_{zi}^+(z)(1 - 2z^{-1} + 3z^{-2}) = 2y(-1) - 3y(-1)z^{-1} - 3y(-2)
\]
Substitute $y(-1) = 1$ and $y(-2) = 1$:
\[
Y_{zi}^+(z)(1 - 2z^{-1} + 3z^{-2}) = 2(1) - 3(1)z^{-1} - 3(1) = -1 - 3z^{-1}
\]
Solve for $Y_{zi}^+(z)$:
\[
Y_{zi}^+(z) = \frac{-1 - 3z^{-1}}{1 - 2z^{-1} + 3z^{-2}}
\]
Using the same parameters $a = \sqrt{3}$, $a\cos\omega_0 = 1$, and $a\sin\omega_0 = \sqrt{2}$ from Part (b), rewrite the numerator as $-1 - 3z^{-1} = -(1 - z^{-1}) - 4z^{-1}$:
\[
\begin{aligned}
Y_{zi}^+(z) &= -\frac{1 - az^{-1}\cos\omega_0}{1 - 2az^{-1}\cos\omega_0 + a^2 z^{-2}} \\
&\quad - 2\sqrt{2}\frac{az^{-1}\sin\omega_0}{1 - 2az^{-1}\cos\omega_0 + a^2 z^{-2}}
\end{aligned}
\]
Invert using Table~\ref{tab:z_pairs} to obtain the zero-input response:
\[
y_{zi}(n) = -(\sqrt{3})^n \left[\cos(\omega_0 n) + 2\sqrt{2}\sin(\omega_0 n)\right]u(n), \quad \omega_0 \approx 54.74^\circ
\]
In single-cosine form, the partial-fraction coefficient at $p_1 = 1 + j\sqrt{2}$ is:
\[
D_1 = \left. \frac{-z - 3}{z - p_2} \right|_{z = 1 + j\sqrt{2}} = -\frac{1}{2} + j\sqrt{2} = \frac{3}{2}e^{j109.47^\circ}
\]
which gives:
\[
y_{zi}(n) = 3(\sqrt{3})^n \cos(54.74^\circ n + 109.47^\circ)u(n)
\]

\subsubsection*{d)}
From Table~\ref{tab:z_pairs}, the $z$-transform of $x(n) = 2^n u(n)$ is:
\[
X(z) = \frac{1}{1 - 2z^{-1}}, \quad \text{ROC: } |z| > 2
\]
The $z$-transform of the zero-state response is the product $Y_{zs}(z) = H(z)X(z)$:
\[
Y_{zs}(z) = \frac{1 + z^{-1}}{(1 - 2z^{-1} + 3z^{-2})(1 - 2z^{-1})}
\]
Multiply the numerator and denominator by $z^3$ and divide both sides by $z$:
\[
\frac{Y_{zs}(z)}{z} = \frac{z(z + 1)}{(z - 2)(z^2 - 2z + 3)} = \frac{Q_1}{z - 2} + \frac{B_1 z + B_2}{z^2 - 2z + 3}
\]
Find the coefficient $Q_1$ for the pole at $z = 2$:
\[
Q_1 = \left. \frac{z(z + 1)}{z^2 - 2z + 3} \right|_{z = 2} = \frac{2(3)}{4 - 4 + 3} = \frac{6}{3} = 2
\]
Subtract $\frac{2}{z - 2}$ from $\frac{Y_{zs}(z)}{z}$ to find the remaining quadratic term:
\[
\begin{aligned}
\frac{Y_{zs}(z)}{z} - \frac{2}{z - 2} &= \frac{z^2 + z - 2(z^2 - 2z + 3)}{(z - 2)(z^2 - 2z + 3)} \\
&= \frac{-z^2 + 5z - 6}{(z - 2)(z^2 - 2z + 3)} \\
&= \frac{-(z - 2)(z - 3)}{(z - 2)(z^2 - 2z + 3)} = \frac{-z + 3}{z^2 - 2z + 3}
\end{aligned}
\]
Multiply both sides by $z$ and rewrite $Y_{zs}(z)$ in powers of $z^{-1}$:
\[
Y_{zs}(z) = \frac{2}{1 - 2z^{-1}} + \frac{-1 + 3z^{-1}}{1 - 2z^{-1} + 3z^{-2}}
\]
Split the numerator of the second fraction as $-1 + 3z^{-1} = -(1 - z^{-1}) + 2z^{-1}$ so it matches the cosine and sine transform pairs with $a = \sqrt{3}$, $a\cos\omega_0 = 1$, and $a\sin\omega_0 = \sqrt{2}$:
\[
\begin{aligned}
Y_{zs}(z) &= \frac{2}{1 - 2z^{-1}} - \frac{1 - az^{-1}\cos\omega_0}{1 - 2az^{-1}\cos\omega_0 + a^2 z^{-2}} \\
&\quad + \sqrt{2}\frac{az^{-1}\sin\omega_0}{1 - 2az^{-1}\cos\omega_0 + a^2 z^{-2}}
\end{aligned}
\]
Invert term by term using Table~\ref{tab:z_pairs} to obtain the causal zero-state response:
\[
y_{zs}(n) = \left[2(2)^n + (\sqrt{3})^n \left(-\cos(\omega_0 n) + \sqrt{2}\sin(\omega_0 n)\right)\right]u(n)
\]
where $\omega_0 = \arccos(1/\sqrt{3}) \approx 54.74^\circ$. In single-cosine form, the coefficient at $p_1 = 1 + j\sqrt{2}$ is:
\[
A_1 = \left. \frac{-z + 3}{z - p_2} \right|_{z = 1 + j\sqrt{2}} = -\frac{1}{2} - j\frac{\sqrt{2}}{2} = -\frac{\sqrt{3}}{2}e^{j\omega_0}
\]
which gives:
\[
y_{zs}(n) = \left[2^{n+1} - (\sqrt{3})^{n+1}\cos\left(\omega_0(n + 1)\right)\right]u(n)
\]

% -------------------------------------------------------------------------
% Exercise 6
% -------------------------------------------------------------------------
\subsection{Exercise 6: Mathematical Proofs of Z-Transform Properties}
\textbf{Problem Statement:} Let $x(n)$ be a discrete-time sequence with $\mathcal{Z}$-transform $X(z) = \sum_{n=-\infty}^{\infty} x(n) z^{-n}$ and Region of Convergence $R_x$. Prove the following analytical statements:
\begin{enumerate}[label=(\alph*)]
    \item $\mathcal{Z}\{x^*(n)\} = X^*(z^*)$
    \item $\mathcal{Z}\{\text{Re}\{x(n)\}\} = \dfrac{1}{2}\left[ X(z) + X^*(z^*) \right]$
    \item $\mathcal{Z}\{\text{Im}\{x(n)\}\} = \dfrac{1}{2j}\left[ X(z) - X^*(z^*) \right]$
    \item $\mathcal{Z}\{e^{j\omega_0 n} x(n)\} = X(z e^{-j\omega_0})$
\end{enumerate}

\subsubsection{Proof of Statement (a): Complex Conjugation}
From the basic definition of the bilateral $\mathcal{Z}$-transform:
\begin{equation}
    \mathcal{Z}\{x^*(n)\} = \sum_{n=-\infty}^{\infty} x^*(n) z^{-n}.
\end{equation}
Recall that for any complex numbers $u$ and $v$, $u^* v^* = (u v)^*$ and $(z^{-n}) = ((z^*)^{-n})^*$. Hence, we can factor out the complex conjugate across the infinite sum:
\begin{equation}
    \mathcal{Z}\{x^*(n)\} = \sum_{n=-\infty}^{\infty} \left[ x(n) (z^*)^{-n} \right]^* = \left( \sum_{n=-\infty}^{\infty} x(n) (z^*)^{-n} \right)^*.
\end{equation}
Recognizing that the summation inside the outer parentheses is precisely the evaluation of $X(\cdot)$ at the conjugate variable $z^*$:
\begin{equation}
    \sum_{n=-\infty}^{\infty} x(n) (z^*)^{-n} = X(z^*).
\end{equation}
Taking the outer complex conjugate gives
\begin{equation}
    \mathcal{Z}\{x^*(n)\} = X^*(z^*).
\end{equation}
The convergence condition requires $z^* \in R_x$. Since $|z^*| = |z|$, the magnitude condition for convergence is identical, so the ROC remains exactly $R_x$. \hfill $\blacksquare$

\subsubsection{Proof of Statement (b): Real Part Extraction}
By Euler's identity, the real part of any complex sequence $x(n)$ is expressed as:
\begin{equation}
    \text{Re}\{x(n)\} = \frac{x(n) + x^*(n)}{2}.
\end{equation}
Applying the linearity property of the $\mathcal{Z}$-transform:
\begin{equation}
    \mathcal{Z}\{\text{Re}\{x(n)\}\} = \mathcal{Z}\left\{ \frac{1}{2} x(n) + \frac{1}{2} x^*(n) \right\} = \frac{1}{2} \mathcal{Z}\{x(n)\} + \frac{1}{2} \mathcal{Z}\{x^*(n)\}.
\end{equation}
Substituting the result of statement (a), $\mathcal{Z}\{x^*(n)\} = X^*(z^*)$:
\begin{equation}
    \mathcal{Z}\{\text{Re}\{x(n)\}\} = \frac{1}{2} \left[ X(z) + X^*(z^*) \right].
\end{equation}
The ROC contains at least $R_x$. \hfill $\blacksquare$

\subsubsection{Proof of Statement (c): Imaginary Part Extraction}
Similarly, the imaginary part of $x(n)$ is defined as:
\begin{equation}
    \text{Im}\{x(n)\} = \frac{x(n) - x^*(n)}{2j}.
\end{equation}
Invoking linearity:
\begin{equation}
    \mathcal{Z}\{\text{Im}\{x(n)\}\} = \frac{1}{2j} \mathcal{Z}\{x(n)\} - \frac{1}{2j} \mathcal{Z}\{x^*(n)\}.
\end{equation}
Substituting $\mathcal{Z}\{x^*(n)\} = X^*(z^*)$ yields:
\begin{equation}
    \mathcal{Z}\{\text{Im}\{x(n)\}\} = \frac{1}{2j} \left[ X(z) - X^*(z^*) \right].
\end{equation}
The ROC contains at least $R_x$. \hfill $\blacksquare$

\subsubsection{Proof of Statement (d): Complex Frequency Modulation}
By direct evaluation using the $\mathcal{Z}$-transform definition:
\begin{equation}
    \mathcal{Z}\{e^{j\omega_0 n} x(n)\} = \sum_{n=-\infty}^{\infty} \left[ e^{j\omega_0 n} x(n) \right] z^{-n}.
\end{equation}
Combining the exponent $n$ terms:
\begin{equation}
    \mathcal{Z}\{e^{j\omega_0 n} x(n)\} = \sum_{n=-\infty}^{\infty} x(n) \left( z e^{-j\omega_0} \right)^{-n}.
\end{equation}
By inspection, this summation is the definition of $X(\cdot)$ evaluated at the scaled complex variable $z' = z e^{-j\omega_0}$:
\begin{equation}
    \mathcal{Z}\{e^{j\omega_0 n} x(n)\} = X\left( z e^{-j\omega_0} \right).
\end{equation}
The Region of Convergence requires $z e^{-j\omega_0} \in R_x$. Because $|z e^{-j\omega_0}| = |z| \cdot |e^{-j\omega_0}| = |z|$, the radial bounds are unchanged, meaning the ROC remains $R_x$. Geometrically, this property corresponds to a counter-clockwise rotation of the pole-zero constellation by the angle $\omega_0$. \hfill $\blacksquare$

\newpage

% =========================================================================
% PART II: ADDITIONAL EXERCISES (PROAKIS & MANOLAKIS CHAPTER 3)
% =========================================================================
\section{Part II: Additional Exercises (Reference Document)}

% -------------------------------------------------------------------------
% Problem 3.3
% -------------------------------------------------------------------------
\subsection{Problem 3.3: Signals and Regions of Convergence}
\textbf{Problem Statement:} Determine the $\mathcal{Z}$-transforms and sketch the Region of Convergence (ROC) for each of the following signals:
\begin{enumerate}[label=(\alph*)]
    \item $x_1(n) = \begin{cases} \left(\dfrac{1}{3}\right)^n, & n \ge 0 \\ \left(\dfrac{1}{2}\right)^{-n}, & n < 0 \end{cases}$
    \item $x_2(n) = \begin{cases} \left(\dfrac{1}{3}\right)^n - 2^n, & n \ge 0 \\ 0, & n < 0 \end{cases}$
    \item $x_3(n) = x_1(n+4)$
    \item $x_4(n) = x_1(-n)$
\end{enumerate}

\subsubsection*{a)}
Since $\left(\frac{1}{2}\right)^{-n} = 2^n$, express $x_1(n)$ as the sum of a causal signal and an anticausal signal:
\[
x_1(n) = \left(\frac{1}{3}\right)^n u(n) + 2^n u(-n - 1)
\]
From Table~\ref{tab:z_pairs}, the $z$-transforms of the two components are:
\[
\left(\frac{1}{3}\right)^n u(n) \stackrel{z}{\longleftrightarrow} \frac{1}{1 - \frac{1}{3}z^{-1}}, \quad \text{ROC: } |z| > \frac{1}{3}
\]
\[
2^n u(-n - 1) \stackrel{z}{\longleftrightarrow} -\frac{1}{1 - 2z^{-1}}, \quad \text{ROC: } |z| < 2
\]
Combine the two transforms and take the intersection of their regions of convergence:
\[
X_1(z) = \frac{1}{1 - \frac{1}{3}z^{-1}} - \frac{1}{1 - 2z^{-1}} = \frac{-\frac{5}{3}z^{-1}}{\left(1 - \frac{1}{3}z^{-1}\right)(1 - 2z^{-1})}
\]
\[
\text{ROC: } \frac{1}{3} < |z| < 2
\]

\begin{figure}[H]
\centering
\begin{tikzpicture}[scale=0.95]
  \fill[pattern=north east lines, pattern color=gray!70, even odd rule]
    (0,0) circle (2.0) (0,0) circle (0.7);
  \draw[dashed, thick] (0,0) circle (2.0);
  \draw[dashed, thick] (0,0) circle (0.7);
  \draw[->] (-2.8,0) -- (2.8,0) node[right] {$\text{Re}(z)$};
  \draw[->] (0,-2.6) -- (0,2.6) node[above] {$\text{Im}(z)$};
  \node[below right] at (0.7,0) {$\frac{1}{3}$};
  \node[below right] at (2.0,0) {$2$};
\end{tikzpicture}
\caption{ROC for $X_1(z)$: $\frac{1}{3} < |z| < 2$.}
\label{fig:p33_roc_a}
\end{figure}

\subsubsection*{b)}
Write $x_2(n)$ using the unit step sequence:
\[
x_2(n) = \left(\frac{1}{3}\right)^n u(n) - 2^n u(n)
\]
By the linearity property in Table~\ref{tab:z_properties} and the pairs in Table~\ref{tab:z_pairs}:
\[
\left(\frac{1}{3}\right)^n u(n) \stackrel{z}{\longleftrightarrow} \frac{1}{1 - \frac{1}{3}z^{-1}}, \quad \text{ROC: } |z| > \frac{1}{3}
\]
\[
2^n u(n) \stackrel{z}{\longleftrightarrow} \frac{1}{1 - 2z^{-1}}, \quad \text{ROC: } |z| > 2
\]
The overall ROC is the intersection of $|z| > \frac{1}{3}$ and $|z| > 2$, which is $|z| > 2$:
\[
X_2(z) = \frac{1}{1 - \frac{1}{3}z^{-1}} - \frac{1}{1 - 2z^{-1}} = \frac{-\frac{5}{3}z^{-1}}{\left(1 - \frac{1}{3}z^{-1}\right)(1 - 2z^{-1})}
\]
\[
\text{ROC: } |z| > 2
\]

\begin{figure}[H]
\centering
\begin{tikzpicture}[scale=0.95]
  \fill[pattern=north east lines, pattern color=gray!70, even odd rule]
    (-2.5,-2.3) rectangle (2.5,2.3) (0,0) circle (1.7);
  \draw[dashed, thick] (0,0) circle (1.7);
  \draw[dashed] (0,0) circle (0.6);
  \draw[->] (-2.8,0) -- (2.8,0) node[right] {$\text{Re}(z)$};
  \draw[->] (0,-2.6) -- (0,2.6) node[above] {$\text{Im}(z)$};
  \node[below right] at (0.6,0) {$\frac{1}{3}$};
  \node[below right] at (1.7,0) {$2$};
\end{tikzpicture}
\caption{ROC for $X_2(z)$: $|z| > 2$.}
\label{fig:p33_roc_b}
\end{figure}

\subsubsection*{c)}
Apply the time-shifting property $x_1(n - k) \stackrel{z}{\longleftrightarrow} z^{-k}X_1(z)$ from Table~\ref{tab:z_properties} with $k = -4$:
\[
X_3(z) = z^4 X_1(z) = \frac{-\frac{5}{3}z^3}{\left(1 - \frac{1}{3}z^{-1}\right)(1 - 2z^{-1})}
\]
Multiplying by $z^4$ can only remove $z = \infty$ from the ROC. Because $z = \infty$ is already outside the ring of $X_1(z)$, the ROC is unchanged:
\[
\text{ROC: } \frac{1}{3} < |z| < 2
\]

\begin{figure}[H]
\centering
\begin{tikzpicture}[scale=0.95]
  \fill[pattern=north east lines, pattern color=gray!70, even odd rule]
    (0,0) circle (2.0) (0,0) circle (0.7);
  \draw[dashed, thick] (0,0) circle (2.0);
  \draw[dashed, thick] (0,0) circle (0.7);
  \draw[->] (-2.8,0) -- (2.8,0) node[right] {$\text{Re}(z)$};
  \draw[->] (0,-2.6) -- (0,2.6) node[above] {$\text{Im}(z)$};
  \node[below right] at (0.7,0) {$\frac{1}{3}$};
  \node[below right] at (2.0,0) {$2$};
\end{tikzpicture}
\caption{ROC for $X_3(z)$: $\frac{1}{3} < |z| < 2$.}
\label{fig:p33_roc_c}
\end{figure}

\subsubsection*{d)}
Apply the time-reversal property $x_1(-n) \stackrel{z}{\longleftrightarrow} X_1(z^{-1})$ from Table~\ref{tab:z_properties}:
\[
\begin{aligned}
X_4(z) &= X_1(z^{-1}) = \frac{1}{1 - \frac{1}{3}z} - \frac{1}{1 - 2z} \\
&= \frac{-\frac{5}{3}z}{\left(1 - \frac{1}{3}z\right)(1 - 2z)} = \frac{-\frac{5}{2}z^{-1}}{\left(1 - \frac{1}{2}z^{-1}\right)(1 - 3z^{-1})}
\end{aligned}
\]
The ROC is found by replacing $z$ with $z^{-1}$ in $\frac{1}{3} < |z| < 2$, which gives $\frac{1}{3} < |z^{-1}| < 2$:
\[
\text{ROC: } \frac{1}{2} < |z| < 3
\]

\begin{figure}[H]
\centering
\begin{tikzpicture}[scale=0.95]
  \fill[pattern=north east lines, pattern color=gray!70, even odd rule]
    (0,0) circle (2.1) (0,0) circle (0.7);
  \draw[dashed, thick] (0,0) circle (2.1);
  \draw[dashed, thick] (0,0) circle (0.7);
  \draw[->] (-2.8,0) -- (2.8,0) node[right] {$\text{Re}(z)$};
  \draw[->] (0,-2.6) -- (0,2.6) node[above] {$\text{Im}(z)$};
  \node[below right] at (0.7,0) {$\frac{1}{2}$};
  \node[below right] at (2.1,0) {$3$};
\end{tikzpicture}
\caption{ROC for $X_4(z)$: $\frac{1}{2} < |z| < 3$.}
\label{fig:p33_roc_d}
\end{figure}

% -------------------------------------------------------------------------
% Problem 3.4
% -------------------------------------------------------------------------
\subsection{Problem 3.4: Determination of Z-Transforms}
\textbf{Problem Statement:} Determine the $\mathcal{Z}$-transform and Region of Convergence for each of the following discrete-time signals:

\subsubsection*{a)}
Given signal:
\begin{equation}
    x(n) = n(-1)^n u(n).
\end{equation}
We know the basic causal exponential pair:
\begin{equation}
    v(n) = (-1)^n u(n) \longleftrightarrow V(z) = \frac{1}{1 + z^{-1}}, \quad \text{ROC: } |z| > 1.
\end{equation}
Using the differentiation property $\mathcal{Z}\{n v(n)\} = -z \frac{d}{dz} V(z)$:
\begin{align}
    X(z) &= -z \frac{d}{dz}\left( \frac{1}{1 + z^{-1}} \right) = -z \left[ -\frac{-z^{-2}}{(1 + z^{-1})^2} \right] \nonumber \\
         &= \frac{-z^{-1}}{(1 + z^{-1})^2} = \frac{-z}{(z + 1)^2}, \quad \text{ROC: } |z| > 1.
\end{align}

\subsubsection*{b)}
Given signal:
\begin{equation}
    x(n) = n^2 u(n).
\end{equation}
We apply the differentiation property successively. For $v(n) = n u(n)$:
\begin{equation}
    V(z) = \mathcal{Z}\{n u(n)\} = \frac{z^{-1}}{(1 - z^{-1})^2}, \quad |z| > 1.
\end{equation}
Now, $x(n) = n \cdot v(n)$, so $X(z) = -z \frac{d}{dz} V(z)$:
\begin{align*}
    \frac{d}{dz} \left[ z^{-1}(1 - z^{-1})^{-2} \right] &= -z^{-2}(1 - z^{-1})^{-2} + z^{-1}\left[ -2(1 - z^{-1})^{-3}(z^{-2}) \right] \\
    &= \frac{-z^{-2}(1 - z^{-1}) - 2z^{-3}}{(1 - z^{-1})^3} = \frac{-z^{-2} - z^{-3}}{(1 - z^{-1})^3}.
\end{align*}
Multiplying by $-z$:
\begin{equation}
    X(z) = \frac{z^{-1} + z^{-2}}{(1 - z^{-1})^3} = \frac{z^{-1}(1 + z^{-1})}{(1 - z^{-1})^3} = \frac{z(z + 1)}{(z - 1)^3}, \quad \text{ROC: } |z| > 1.
\end{equation}

\subsubsection*{c)}
Given signal:
\begin{equation}
    x(n) = -n a^n u(-n-1).
\end{equation}
For the anti-causal exponential sequence:
\begin{equation}
    v(n) = -a^n u(-n-1) \longleftrightarrow V(z) = \frac{1}{1 - a z^{-1}}, \quad \text{ROC: } |z| < |a|.
\end{equation}
Because $x(n) = n v(n)$, we apply the differentiation property:
\begin{align}
    X(z) &= -z \frac{d}{dz}\left( \frac{1}{1 - a z^{-1}} \right) = -z \left[ -\frac{a z^{-2}}{(1 - a z^{-1})^2} \right] \nonumber \\
         &= \frac{a z^{-1}}{(1 - a z^{-1})^2} = \frac{a z}{(z - a)^2}, \quad \text{ROC: } |z| < |a|.
\end{align}

\subsubsection*{d)}
Given signal:
\begin{equation}
    x(n) = (-1)^n \left(\cos \frac{\pi}{3} n\right) u(n).
\end{equation}
Recall the standard cosine transform pair:
\begin{equation}
    \cos(\omega_0 n) u(n) \longleftrightarrow \frac{1 - (\cos\omega_0) z^{-1}}{1 - 2(\cos\omega_0) z^{-1} + z^{-2}}, \quad |z| > 1.
\end{equation}
For $\omega_0 = \frac{\pi}{3}$, $\cos \frac{\pi}{3} = \frac{1}{2}$. Thus:
\begin{equation}
    V(z) = \mathcal{Z}\left\{ \cos\left(\frac{\pi}{3} n\right) u(n) \right\} = \frac{1 - \frac{1}{2}z^{-1}}{1 - z^{-1} + z^{-2}}, \quad |z| > 1.
\end{equation}
Applying the scaling property in the $\mathcal{Z}$-domain, $\mathcal{Z}\{a^n v(n)\} = V(a^{-1}z)$. Here $a = -1$, which corresponds to replacing $z^{-1}$ by $-z^{-1}$:
\begin{equation}
    X(z) = V(-z) = \frac{1 - \frac{1}{2}(-z^{-1})}{1 - (-z^{-1}) + (-z^{-1})^2} = \frac{1 + \frac{1}{2}z^{-1}}{1 + z^{-1} + z^{-2}}, \quad \text{ROC: } |z| > 1.
\end{equation}

\subsubsection*{e)}
Given signal:
\begin{equation}
    x(n) = (-1)^n u(n).
\end{equation}
Evaluating the geometric series:
\begin{equation}
    X(z) = \sum_{n=0}^{\infty} (-1)^n z^{-n} = \sum_{n=0}^{\infty} (-z^{-1})^n = \frac{1}{1 + z^{-1}} = \frac{z}{z + 1}, \quad \text{ROC: } |z| > 1.
\end{equation}

\subsubsection*{f)}
Given signal:
\begin{equation}
    x(n) = \{\org{1}, 0, -1, 0, 1, -1, \dots\}.
\end{equation}
Depending on the interpretation of the trailing terms:
\begin{enumerate}
    \item \textbf{Finite-Duration Sequence interpretation:} If $x(n)$ denotes the specified 6-sample sequence $x(n) = \{\org{1}, 0, -1, 0, 1, -1\}$:
    \begin{equation}
        X(z) = 1 - z^{-2} + z^{-4} - z^{-5}, \quad \text{ROC: } 0 < |z| < \infty.
    \end{equation}
    \item \textbf{Causal Sinusoidal Sequence interpretation:} If the series represents the alternating trigonometric sequence $x(n) = \cos\left(\frac{\pi}{2} n\right) u(n)$:
    \begin{equation}
        x(n) = \{\org{1}, 0, -1, 0, 1, 0, -1, \dots\} \longleftrightarrow X(z) = \frac{1}{1 + z^{-2}} = \frac{z^2}{z^2 + 1}, \quad \text{ROC: } |z| > 1.
    \end{equation}
    \item \textbf{Periodic Sequence with Period $N=6$:} If $x(n)$ repeats periodically with fundamental block $[1, 0, -1, 0, 1, -1]$ for $n \ge 0$:
    \begin{equation}
        X(z) = \frac{1 - z^{-2} + z^{-4} - z^{-5}}{1 - z^{-6}}, \quad \text{ROC: } |z| > 1.
    \end{equation}
\end{enumerate}

% -------------------------------------------------------------------------
% Problem 3.7
% -------------------------------------------------------------------------
\subsection{Problem 3.7: Convolution by Means of Z-Transform}
\textbf{Problem Statement:} Compute the convolution $y(n) = x_1(n) * x_2(n)$ by means of the $\mathcal{Z}$-transform for the signals:
\begin{equation}
    x_1(n) = \begin{cases} \left(\dfrac{1}{3}\right)^n, & n \ge 0 \\ \left(\dfrac{1}{2}\right)^{-n}, & n < 0 \end{cases}, \qquad x_2(n) = \left(\frac{1}{2}\right)^n u(n).
\end{equation}

\textbf{Solution:}
From Problem 3.3(a), the $z$-transform of $x_1(n)$ is:
\[
X_1(z) = \frac{-\frac{5}{3}z^{-1}}{\left(1 - \frac{1}{3}z^{-1}\right)(1 - 2z^{-1})}, \quad \text{ROC: } \frac{1}{3} < |z| < 2
\]
From Table~\ref{tab:z_pairs}, the $z$-transform of $x_2(n)$ is:
\[
X_2(z) = \frac{1}{1 - \frac{1}{2}z^{-1}}, \quad \text{ROC: } |z| > \frac{1}{2}
\]
By the convolution property in Table~\ref{tab:z_properties}, $Y(z) = X_1(z)X_2(z)$:
\[
Y(z) = \frac{-\frac{5}{3}z^{-1}}{\left(1 - \frac{1}{3}z^{-1}\right)\left(1 - \frac{1}{2}z^{-1}\right)(1 - 2z^{-1})}
\]
The ROC of $Y(z)$ is the intersection of $\frac{1}{3} < |z| < 2$ and $|z| > \frac{1}{2}$:
\[
\text{ROC: } \frac{1}{2} < |z| < 2
\]
Multiply the numerator and denominator of $Y(z)$ by $z^3$ and divide both sides by $z$:
\[
\begin{aligned}
\frac{Y(z)}{z} &= \frac{-\frac{5}{3}z}{\left(z - \frac{1}{3}\right)\left(z - \frac{1}{2}\right)(z - 2)} \\
&= \frac{A_1}{z - \frac{1}{3}} + \frac{A_2}{z - \frac{1}{2}} + \frac{A_3}{z - 2}
\end{aligned}
\]
Evaluate the partial-fraction coefficients at the three poles:
\[
A_1 = \left. \frac{-\frac{5}{3}z}{\left(z - \frac{1}{2}\right)(z - 2)} \right|_{z = 1/3} = \frac{-\frac{5}{9}}{\left(-\frac{1}{6}\right)\left(-\frac{5}{3}\right)} = -2
\]
\[
A_2 = \left. \frac{-\frac{5}{3}z}{\left(z - \frac{1}{3}\right)(z - 2)} \right|_{z = 1/2} = \frac{-\frac{5}{6}}{\left(\frac{1}{6}\right)\left(-\frac{3}{2}\right)} = \frac{10}{3}
\]
\[
A_3 = \left. \frac{-\frac{5}{3}z}{\left(z - \frac{1}{3}\right)\left(z - \frac{1}{2}\right)} \right|_{z = 2} = \frac{-\frac{10}{3}}{\left(\frac{5}{3}\right)\left(\frac{3}{2}\right)} = -\frac{4}{3}
\]
Multiply by $z$ and write $Y(z)$ in powers of $z^{-1}$:
\[
Y(z) = -2\frac{1}{1 - \frac{1}{3}z^{-1}} + \frac{10}{3}\frac{1}{1 - \frac{1}{2}z^{-1}} - \frac{4}{3}\frac{1}{1 - 2z^{-1}}
\]
Since the ROC is $\frac{1}{2} < |z| < 2$, it lies outside the poles at $\frac{1}{3}$ and $\frac{1}{2}$ (which produce causal terms) and inside the pole at $2$ (which produces an anticausal term):
\[
y(n) = \left[-2\left(\frac{1}{3}\right)^n + \frac{10}{3}\left(\frac{1}{2}\right)^n\right]u(n) + \frac{4}{3}(2)^n u(-n - 1)
\]

% -------------------------------------------------------------------------
% Problem 3.8
% -------------------------------------------------------------------------
\subsection{Problem 3.8: Accumulation and Convolution Properties}
\textbf{Problem Statement:}
\begin{enumerate}[label=(\alph*)]
    \item Use the convolution property to express the $\mathcal{Z}$-transform of the accumulation signal
    \begin{equation}
        y(n) = \sum_{k=-\infty}^{n} x(k)
        \label{eq:p38_accum}
    \end{equation}
    in terms of $X(z)$.
    \item Determine the $\mathcal{Z}$-transform of $x(n) = (n+1)u(n)$ using the convolution property. [Hint: Show first that $x(n) = u(n) * u(n)$.]
\end{enumerate}

\subsubsection*{a)}
The accumulation operation in Eq.~\eqref{eq:p38_accum} can be written in convolution form:
\begin{equation}
    y(n) = \sum_{k=-\infty}^{\infty} x(k) u(n - k) = x(n) * u(n).
\end{equation}
By the convolution theorem of the $\mathcal{Z}$-transform:
\begin{equation}
    Y(z) = \mathcal{Z}\{x(n) * u(n)\} = X(z) U(z).
\end{equation}
Since $U(z) = \mathcal{Z}\{u(n)\} = \frac{1}{1 - z^{-1}}$ with $\text{ROC}_u: |z| > 1$:
\begin{equation}
    Y(z) = \frac{X(z)}{1 - z^{-1}} = \frac{z}{z - 1} X(z).
\end{equation}
The Region of Convergence is at least $\text{ROC}_x \cap \{z \in \mathbb{C} : |z| > 1\}$, with possible extension if $X(z)$ has a zero that cancels the pole at $z = 1$.

\subsubsection*{b)}
\textbf{Step 1: Verification of the Hint.} Compute the direct discrete convolution of $u(n)$ with itself:
\begin{equation}
    (u * u)(n) = \sum_{k=-\infty}^{\infty} u(k) u(n - k).
\end{equation}
The product $u(k)u(n-k)$ is non-zero only when both $k \ge 0$ and $n - k \ge 0 \implies 0 \le k \le n$.
\begin{itemize}
    \item For $n < 0$: No integers $k$ satisfy $0 \le k \le n$, hence $(u * u)(n) = 0$.
    \item For $n \ge 0$: The sum runs from $k = 0$ to $k = n$:
    \begin{equation}
        (u * u)(n) = \sum_{k=0}^{n} 1 = n + 1.
    \end{equation}
\end{itemize}
Combining the cases confirms:
\begin{equation}
    u(n) * u(n) = (n + 1)u(n) = x(n). \hfill \blacksquare
\end{equation}

\textbf{Step 2: Transform Calculation.}
Using the convolution theorem:
\begin{equation}
    X(z) = \mathcal{Z}\{u(n) * u(n)\} = U(z) \cdot U(z) = [U(z)]^2.
\end{equation}
Substituting $U(z) = \frac{1}{1 - z^{-1}}$ for $|z| > 1$:
\begin{equation}
    X(z) = \left( \frac{1}{1 - z^{-1}} \right)^2 = \frac{1}{(1 - z^{-1})^2} = \frac{z^2}{(z - 1)^2}, \quad \text{ROC: } |z| > 1.
\end{equation}
This matches the derivative property result: $\mathcal{Z}\{(n+1)u(n)\} = \mathcal{Z}\{n u(n)\} + \mathcal{Z}\{u(n)\} = \frac{z^{-1}}{(1 - z^{-1})^2} + \frac{1}{1 - z^{-1}} = \frac{1}{(1 - z^{-1})^2}$.

% -------------------------------------------------------------------------
% Problem 3.13
% -------------------------------------------------------------------------
\subsection{Problem 3.13: Upsampling and Downsampling Properties}
\textbf{Problem Statement:} Let $x(n)$ be a sequence with $\mathcal{Z}$-transform $X(z)$. Determine, in terms of $X(z)$, the $\mathcal{Z}$-transforms of the following signals:
\begin{enumerate}[label=(\alph*)]
    \item $x_1(n) = \begin{cases} x\left(\dfrac{n}{2}\right), & \text{if } n \text{ is even} \\ 0, & \text{if } n \text{ is odd} \end{cases}$ \quad (Upsampling by factor 2)
    \item $x_2(n) = x(2n)$ \quad (Downsampling / Decimation by factor 2)
\end{enumerate}

\subsubsection*{a)}
Write the definition of the $z$-transform for $x_1(n)$:
\[
X_1(z) = \sum_{n=-\infty}^{\infty} x_1(n)z^{-n}
\]
Since $x_1(n) = 0$ for all odd $n$, substitute $n = 2k$, where $k$ runs through all integers:
\[
X_1(z) = \sum_{k=-\infty}^{\infty} x(k)z^{-2k} = \sum_{k=-\infty}^{\infty} x(k)\left(z^2\right)^{-k} = X\left(z^2\right)
\]

\subsubsection*{b)}
Write the definition of the $z$-transform for $x_2(n)$:
\[
X_2(z) = \sum_{n=-\infty}^{\infty} x_2(n)z^{-n} = \sum_{n=-\infty}^{\infty} x(2n)z^{-n}
\]
Let $k = 2n$. The factor $\frac{1}{2}\left[1 + (-1)^k\right]$ equals $1$ when $k$ is even and $0$ when $k$ is odd. Using this factor, extend the summation over all integers $k$:
\[
X_2(z) = \sum_{k=-\infty}^{\infty} \frac{1}{2}\left[1 + (-1)^k\right]x(k)z^{-k/2}
\]
Separate the sum into two parts:
\[
\begin{aligned}
X_2(z) &= \frac{1}{2}\sum_{k=-\infty}^{\infty} x(k)\left(z^{1/2}\right)^{-k} + \frac{1}{2}\sum_{k=-\infty}^{\infty} x(k)\left(-z^{1/2}\right)^{-k} \\
&= \frac{1}{2}\left[X\left(z^{1/2}\right) + X\left(-z^{1/2}\right)\right]
\end{aligned}
\]

% -------------------------------------------------------------------------
% Problem 3.14
% -------------------------------------------------------------------------
\subsection{Problem 3.14: Causal Inverse Z-Transforms}
\textbf{Problem Statement:} Determine the causal signal $x(n)$ (where the Region of Convergence is exterior to the outermost pole, i.e., $|z| > \max_k |p_k|$) for each of the following $\mathcal{Z}$-transforms:

\subsubsection*{a)}
Given:
\begin{equation}
    X(z) = \frac{1 + 3z^{-1}}{1 + 3z^{-1} + 2z^{-2}}.
\end{equation}
Factoring the denominator polynomial:
\begin{equation}
    1 + 3z^{-1} + 2z^{-2} = (1 + z^{-1})(1 + 2z^{-1}).
\end{equation}
The poles are located at $z = -1$ and $z = -2$. For a causal sequence, the ROC is $|z| > 2$.
Applying partial fraction expansion:
\begin{equation}
    X(z) = \frac{1 + 3z^{-1}}{(1 + z^{-1})(1 + 2z^{-1})} = \frac{A_1}{1 + z^{-1}} + \frac{A_2}{1 + 2z^{-1}}.
\end{equation}
Computing the coefficients:
\begin{align*}
    A_1 &= \left. \frac{1 + 3z^{-1}}{1 + 2z^{-1}} \right|_{z^{-1} = -1} = \frac{1 - 3}{1 - 2} = 2, \\
    A_2 &= \left. \frac{1 + 3z^{-1}}{1 + z^{-1}} \right|_{z^{-1} = -1/2} = \frac{1 - 1.5}{1 - 0.5} = -1.
\end{align*}
Inverting each term for $|z| > 2$:
\begin{equation}
    x(n) = \left[ 2(-1)^n - (-2)^n \right] u(n).
\end{equation}

\subsubsection*{b)}
Given:
\begin{equation}
    X(z) = \frac{1}{1 - z^{-1} + \frac{1}{2}z^{-2}}.
\end{equation}
The denominator is of the standard form $1 - 2r \cos\omega_0 z^{-1} + r^2 z^{-2}$. Comparing coefficients:
\begin{align*}
    r^2 &= \frac{1}{2} \implies r = \frac{1}{\sqrt{2}}, \\
    2r \cos\omega_0 &= 1 \implies 2\left(\frac{1}{\sqrt{2}}\right)\cos\omega_0 = 1 \implies \cos\omega_0 = \frac{1}{\sqrt{2}} \implies \omega_0 = \frac{\pi}{4}.
\end{align*}
Consequently, $r \sin\omega_0 = \frac{1}{\sqrt{2}} \sin\frac{\pi}{4} = \frac{1}{2}$. We decompose the numerator $1$:
\begin{equation}
    1 = \left( 1 - \frac{1}{2}z^{-1} \right) + \frac{1}{2}z^{-1} = \left( 1 - r \cos\omega_0 z^{-1} \right) + \left( r \sin\omega_0 z^{-1} \right).
\end{equation}
Rewriting $X(z)$:
\begin{equation}
    X(z) = \frac{1 - \frac{1}{2}z^{-1}}{1 - z^{-1} + \frac{1}{2}z^{-2}} + \frac{\frac{1}{2}z^{-1}}{1 - z^{-1} + \frac{1}{2}z^{-2}}.
\end{equation}
Taking the causal inverse $\mathcal{Z}$-transform term by term:
\begin{equation}
    x(n) = \left(\frac{1}{\sqrt{2}}\right)^n \left[ \cos\left(\frac{\pi}{4} n\right) + \sin\left(\frac{\pi}{4} n\right) \right] u(n).
\end{equation}
Using the trigonometric identity $\cos\theta + \sin\theta = \sqrt{2} \sin(\theta + \pi/4)$:
\begin{equation}
    x(n) = \sqrt{2}\left(\frac{1}{\sqrt{2}}\right)^n \sin\left(\frac{\pi}{4}(n+1)\right) u(n) = \left(\frac{1}{\sqrt{2}}\right)^{n-1} \sin\left(\frac{\pi}{4}(n+1)\right) u(n).
\end{equation}

\subsubsection*{c)}
Given:
\begin{equation}
    X(z) = \frac{z^{-6} + z^{-7}}{1 - z^{-1}}.
\end{equation}
Expressing $X(z)$ as a product with the unit step transform:
\begin{equation}
    X(z) = (z^{-6} + z^{-7}) \cdot \frac{1}{1 - z^{-1}}.
\end{equation}
Since $\mathcal{Z}^{-1}\left\{ \frac{1}{1 - z^{-1}} \right\} = u(n)$, the time-shifting property gives:
\begin{equation}
    x(n) = u(n - 6) + u(n - 7).
\end{equation}

\subsubsection*{d)}
Given:
\begin{equation}
    X(z) = \frac{1 + 2z^{-2}}{1 + z^{-2}}.
\end{equation}
Performing polynomial division:
\begin{equation}
    X(z) = \frac{2(1 + z^{-2}) - 1}{1 + z^{-2}} = 2 - \frac{1}{1 + z^{-2}}.
\end{equation}
Recall that $\frac{1}{1 + z^{-2}} = \frac{1 - \cos(\pi/2) z^{-1}}{1 - 2\cos(\pi/2) z^{-1} + z^{-2}} \longleftrightarrow \cos\left(\frac{\pi}{2} n\right) u(n)$. Hence:
\begin{equation}
    x(n) = 2\delta(n) - \cos\left(\frac{\pi}{2} n\right) u(n).
\end{equation}

\subsubsection*{e)}
Given:
\begin{equation}
    X(z) = \frac{1}{4} \frac{1 + 6z^{-1} + z^{-2}}{(1 - 2z^{-1} + 2z^{-2})(1 - 0.5z^{-1})}.
\end{equation}
We perform a partial fraction expansion:
\begin{equation}
    X(z) = \frac{A_1 + A_2 z^{-1}}{1 - 2z^{-1} + 2z^{-2}} + \frac{B}{1 - 0.5z^{-1}}.
\end{equation}
Computing $B$ at $z^{-1} = 2$:
\begin{equation}
    B = \left. \frac{1}{4} \frac{1 + 6z^{-1} + z^{-2}}{1 - 2z^{-1} + 2z^{-2}} \right|_{z^{-1} = 2} = \frac{1}{4} \frac{1 + 12 + 4}{1 - 4 + 8} = \frac{17}{20} = 0.85.
\end{equation}
Matching polynomial coefficients on the numerator:
\begin{equation*}
    \frac{1}{4}(1 + 6z^{-1} + z^{-2}) - 0.85(1 - 2z^{-1} + 2z^{-2}) = (A_1 + A_2 z^{-1})(1 - 0.5z^{-1}).
\end{equation*}
Solving gives $A_1 = -0.6$ and $A_2 = 2.9$. For the quadratic denominator $1 - 2z^{-1} + 2z^{-2}$:
\begin{equation*}
    r^2 = 2 \implies r = \sqrt{2}, \quad 2r\cos\omega_0 = 2 \implies \cos\omega_0 = \frac{1}{\sqrt{2}} \implies \omega_0 = \frac{\pi}{4}.
\end{equation*}
Decomposing $-0.6 + 2.9z^{-1} = -0.6(1 - z^{-1}) + 2.3 z^{-1}$:
\begin{align}
    x(n) &= \left[ -0.6 (\sqrt{2})^n \cos\left(\frac{\pi}{4} n\right) + 2.3 (\sqrt{2})^n \sin\left(\frac{\pi}{4} n\right) + 0.85 (0.5)^n \right] u(n) \nonumber \\
         &= \left[ (\sqrt{2})^n \left( -0.6 \cos\left(\frac{\pi}{4} n\right) + 2.3 \sin\left(\frac{\pi}{4} n\right) \right) + \frac{17}{20}\left(\frac{1}{2}\right)^n \right] u(n).
\end{align}

\subsubsection*{f)}
Given:
\begin{equation}
    X(z) = \frac{2 - 1.5z^{-1}}{1 - 1.5z^{-1} + 0.5z^{-2}}.
\end{equation}
Factoring the denominator:
\begin{equation}
    1 - 1.5z^{-1} + 0.5z^{-2} = (1 - z^{-1})(1 - 0.5z^{-1}).
\end{equation}
Expanding by partial fractions:
\begin{equation}
    X(z) = \frac{A}{1 - z^{-1}} + \frac{B}{1 - 0.5z^{-1}}.
\end{equation}
\begin{align*}
    A &= \left. \frac{2 - 1.5z^{-1}}{1 - 0.5z^{-1}} \right|_{z^{-1} = 1} = \frac{0.5}{0.5} = 1, \\
    B &= \left. \frac{2 - 1.5z^{-1}}{1 - z^{-1}} \right|_{z^{-1} = 2} = \frac{2 - 3}{1 - 2} = 1.
\end{align*}
Taking the causal inverse $\mathcal{Z}$-transform:
\begin{equation}
    x(n) = \left[ 1 + (0.5)^n \right] u(n).
\end{equation}

\subsubsection*{g)}
Given:
\begin{equation}
    X(z) = \frac{1 + 2z^{-1} + z^{-2}}{1 + 4z^{-1} + 4z^{-2}}.
\end{equation}
Notice that numerator and denominator are perfect squares: $\frac{(1 + z^{-1})^2}{(1 + 2z^{-1})^2}$. Dividing:
\begin{equation}
    X(z) = \frac{1}{4} + \frac{1}{4} \frac{3 + 4z^{-1}}{(1 + 2z^{-1})^2} = \frac{1}{4} + \frac{1}{2} \frac{1}{1 + 2z^{-1}} + \frac{1}{4} \frac{1}{(1 + 2z^{-1})^2}.
\end{equation}
Using $\mathcal{Z}^{-1}\left\{ \frac{1}{(1 - a z^{-1})^2} \right\} = (n+1) a^n u(n)$ with $a = -2$:
\begin{align}
    x(n) &= \frac{1}{4}\delta(n) + \left[ \frac{1}{2}(-2)^n + \frac{1}{4}(n+1)(-2)^n \right] u(n) \nonumber \\
         &= \frac{1}{4}\delta(n) + \frac{n+3}{4}(-2)^n u(n).
\end{align}

\subsubsection*{h)}
From Figure P3.14 in the reference document:
\begin{itemize}
    \item System gain: $G = \frac{1}{4}$.
    \item Zeros: $z_0 = 0$, $z_1 = -1/2$, and $z_2 = -1/4$.
    \item Poles: $p_1 = 1/2$, and a complex conjugate pair at $p_{2,3} = \frac{1}{\sqrt{2}} e^{\pm j\pi/4}$.
\end{itemize}
The transfer function in terms of $z^{-1}$ is:
\begin{equation}
    X(z) = \frac{1}{4} \frac{(1 + \frac{1}{2}z^{-1})(1 + \frac{1}{4}z^{-1})}{(1 - \frac{1}{2}z^{-1})(1 - z^{-1} + \frac{1}{2}z^{-2})} = \frac{1}{4} \frac{1 + \frac{3}{4}z^{-1} + \frac{1}{8}z^{-2}}{(1 - \frac{1}{2}z^{-1})(1 - z^{-1} + \frac{1}{2}z^{-2})}.
\end{equation}
Partial fraction expansion yields:
\begin{equation}
    X(z) = \frac{3/4}{1 - \frac{1}{2}z^{-1}} + \frac{-0.5 + \frac{11}{16}z^{-1}}{1 - z^{-1} + \frac{1}{2}z^{-2}}.
\end{equation}
Decomposing the second numerator as $-0.5(1 - \frac{1}{2}z^{-1}) + \frac{7}{8}(\frac{1}{2}z^{-1})$:
\begin{equation}
    x(n) = \left[ \frac{3}{4}\left(\frac{1}{2}\right)^n - \frac{1}{2}\left(\frac{1}{\sqrt{2}}\right)^n \cos\left(\frac{\pi}{4} n\right) + \frac{7}{8}\left(\frac{1}{\sqrt{2}}\right)^n \sin\left(\frac{\pi}{4} n\right) \right] u(n).
\end{equation}

\subsubsection*{i)}
Given:
\begin{equation}
    X(z) = \frac{1 - \frac{1}{2}z^{-1}}{1 + \frac{1}{2}z^{-1}}.
\end{equation}
Dividing the numerator by the denominator:
\begin{equation}
    X(z) = -1 + \frac{2}{1 + \frac{1}{2}z^{-1}}.
\end{equation}
Inverting term by term:
\begin{equation}
    x(n) = -\delta(n) + 2\left(-\frac{1}{2}\right)^n u(n).
\end{equation}

\subsubsection*{j)}
Given:
\begin{equation}
    X(z) = \frac{1 - a z^{-1}}{z^{-1} - a}.
\end{equation}
Rewriting in terms of standard causal poles:
\begin{equation}
    X(z) = \frac{1 - a z^{-1}}{-a (1 - a^{-1} z^{-1})} = -a - \frac{1 - a^2}{a (1 - a^{-1} z^{-1})}.
\end{equation}
For a causal sequence, the ROC is exterior to the pole at $z = 1/a$, i.e., $|z| > |1/a|$:
\begin{equation}
    x(n) = -a \delta(n) - \frac{1 - a^2}{a} \left(\frac{1}{a}\right)^n u(n) = -a \delta(n) + \left(a - \frac{1}{a}\right) a^{-n} u(n).
\end{equation}

% -------------------------------------------------------------------------
% Problem 3.19
% -------------------------------------------------------------------------
\subsection{Problem 3.19: Inverse Transform via Differentiation}
\textbf{Problem Statement:} By first differentiating $X(z)$ and then using appropriate properties of the $\mathcal{Z}$-transform, determine the sequence $x(n)$ for each of the following transforms:
\begin{enumerate}[label=(\alph*)]
    \item $X(z) = \log(1 - 2z), \quad |z| < \dfrac{1}{2}$
    \item $X(z) = \log(1 - z^{-1}), \quad |z| > 1$
\end{enumerate}

\subsubsection*{a)}
Differentiate $X(z)$ with respect to $z$:
\[
\frac{dX(z)}{dz} = \frac{-2}{1 - 2z}
\]
Multiply by $-z$ and rewrite in powers of $z^{-1}$:
\[
-z\frac{dX(z)}{dz} = \frac{2z}{1 - 2z} = \frac{2z}{-2z\left(1 - \frac{1}{2}z^{-1}\right)} = -\frac{1}{1 - \frac{1}{2}z^{-1}}, \quad |z| < \frac{1}{2}
\]
By the differentiation property in the $z$-domain from Table~\ref{tab:z_properties}, $nx(n) \stackrel{z}{\longleftrightarrow} -z\frac{dX(z)}{dz}$. Using Table~\ref{tab:z_pairs} for the anticausal region $|z| < \frac{1}{2}$, invert the right-hand side:
\[
nx(n) = \left(\frac{1}{2}\right)^n u(-n - 1)
\]
(At $n = 0$, the power series expansion of $\log(1 - 2z)$ around $z = 0$ has no constant term, so $x(0) = 0$.) Divide by $n$ to obtain $x(n)$:
\[
x(n) = \frac{1}{n}\left(\frac{1}{2}\right)^n u(-n - 1)
\]

\subsubsection*{b)}
Differentiate $X(z)$ with respect to $z$:
\[
\frac{dX(z)}{dz} = \frac{z^{-2}}{1 - z^{-1}}
\]
Multiply both sides by $-z$:
\[
-z\frac{dX(z)}{dz} = -z^{-1}\left(\frac{1}{1 - z^{-1}}\right)
\]
For the exterior causal region ($|z| > 1$), the inverse $z$-transform of $\frac{1}{1 - z^{-1}}$ from Table~\ref{tab:z_pairs} is $u(n)$. By the time-shifting property in Table~\ref{tab:z_properties}, multiplication by $z^{-1}$ delays the sequence by one sample, and by the differentiation property in the $z$-domain:
\[
nx(n) = -u(n - 1)
\]
Divide both sides by $n$ to obtain $x(n)$:
\[
x(n) = -\frac{1}{n}u(n - 1)
\]

% -------------------------------------------------------------------------
% Problem 3.20
% -------------------------------------------------------------------------
\subsection{Problem 3.20: Pole-Zero Pattern Analysis and Synthesis}
\textbf{Problem Statement:}
\begin{enumerate}[label=(\alph*)]
    \item Draw the pole-zero pattern for the signal
    \begin{equation}
        x_1(n) = (r^n \sin\omega_0 n) u(n), \quad 0 < r < 1.
        \label{eq:p320_sig}
    \end{equation}
    \item Compute the $\mathcal{Z}$-transform $X_2(z)$ corresponding directly to the pole-zero pattern in part (a).
    \item Compare $X_1(z)$ with $X_2(z)$. Are they identical? If not, indicate a method to derive $X_1(z)$ from the pole-zero pattern.
\end{enumerate}

\subsubsection*{a)}
Expressing the sinusoidal factor using Euler's formula:
\begin{equation}
    x_1(n) = r^n \left( \frac{e^{j\omega_0 n} - e^{-j\omega_0 n}}{2j} \right) u(n) = \frac{1}{2j}(r e^{j\omega_0})^n u(n) - \frac{1}{2j}(r e^{-j\omega_0})^n u(n).
\end{equation}
Taking the $\mathcal{Z}$-transform of each term:
\begin{align}
    X_1(z) &= \frac{1}{2j} \left[ \frac{1}{1 - r e^{j\omega_0} z^{-1}} - \frac{1}{1 - r e^{-j\omega_0} z^{-1}} \right] \nonumber \\
           &= \frac{1}{2j} \frac{(1 - r e^{-j\omega_0} z^{-1}) - (1 - r e^{j\omega_0} z^{-1})}{(1 - r e^{j\omega_0} z^{-1})(1 - r e^{-j\omega_0} z^{-1})} \nonumber \\
           &= \frac{r \sin\omega_0 z^{-1}}{1 - 2r \cos\omega_0 z^{-1} + r^2 z^{-2}}, \quad \text{ROC: } |z| > r.
\end{align}
Expressed in positive powers of $z$:
\begin{equation}
    X_1(z) = \frac{(r \sin\omega_0) z}{z^2 - 2r \cos\omega_0 z + r^2} = \frac{(r \sin\omega_0) z}{(z - r e^{j\omega_0})(z - r e^{-j\omega_0})}.
\end{equation}
The pole-zero distribution is therefore:
\begin{itemize}
    \item \textbf{Poles:} A complex-conjugate pair at $p_{1,2} = r e^{\pm j\omega_0}$ (located strictly inside the unit circle because $0 < r < 1$).
    \item \textbf{Zeros:} A single finite zero at the origin $z_1 = 0$.
\end{itemize}
The pole-zero constellation in the complex $z$-plane is shown in Figure~\ref{fig:p320_pole_zero}.

\begin{figure}[H]
    \centering
    \begin{tikzpicture}[scale=2.2, >=latex', thick]
        % Axes
        \draw [->] (-1.5,0) -- (1.6,0) node [right] {$\text{Re}(z)$};
        \draw [->] (0,-1.5) -- (0,1.6) node [above] {$\text{Im}(z)$};
        
        % Unit circle
        \draw [dashed, gray!80] (0,0) circle (1.0);
        \node [below right, gray!80] at (1,0) {$1$};
        \node [above left, gray!80] at (0,1) {$j$};
        
        % Pole radius r circle (dotted)
        \draw [dotted, blue!60] (0,0) circle (0.7);
        
        % Pole locations: r=0.7, omega_0 = 45 degrees
        \coordinate (p1) at (45:0.7);
        \coordinate (p2) at (-45:0.7);
        
        % Radial line to pole
        \draw [thin, blue!60] (0,0) -- (p1);
        \draw [thin, blue!60] (0,0) -- (p2);
        
        % Angle arc
        \draw [->, thin] (0.3,0) arc (0:45:0.3);
        \node at (22.5:0.42) {\footnotesize $\omega_0$};
        \draw [->, thin] (0.3,0) arc (0:-45:0.3);
        \node at (-22.5:0.42) {\footnotesize $-\omega_0$};
        
        % Radius label
        \node [above left, blue!80] at (45:0.35) {\footnotesize $r$};

        % Zero at origin
        \draw [line width=1.2pt] (0,0) circle (0.05);
        \node [below left] at (-0.05,-0.05) {\footnotesize $z = 0$};
        
        % Poles (crosses)
        \draw [red, line width=1.4pt] ($(p1)-(0.06,0.06)$) -- ($(p1)+(0.06,0.06)$);
        \draw [red, line width=1.4pt] ($(p1)-(0.06,-0.06)$) -- ($(p1)+(0.06,-0.06)$);
        \node [above right, red] at (p1) {$p_1 = r e^{j\omega_0}$};
        
        \draw [red, line width=1.4pt] ($(p2)-(0.06,0.06)$) -- ($(p2)+(0.06,0.06)$);
        \draw [red, line width=1.4pt] ($(p2)-(0.06,-0.06)$) -- ($(p2)+(0.06,-0.06)$);
        \node [below right, red] at (p2) {$p_2 = r e^{-j\omega_0}$};
    \end{tikzpicture}
    \caption{Pole-zero pattern for $x_1(n) = (r^n \sin\omega_0 n)u(n)$ with $0 < r < 1$.}
    \label{fig:p320_pole_zero}
\end{figure}

\subsubsection*{b)}
From the graphical pole-zero constellation, any rational function with a zero at $z=0$ and conjugate poles at $z = r e^{\pm j\omega_0}$ is constructed in monic rational form as:
\begin{equation}
    X_2(z) = K \frac{z - 0}{(z - r e^{j\omega_0})(z - r e^{-j\omega_0})} = K \frac{z}{z^2 - 2r \cos\omega_0 z + r^2},
\end{equation}
where $K$ represents an arbitrary real or complex scaling constant. Assuming normalized unity gain ($K = 1$):
\begin{equation}
    X_2(z) = \frac{z}{z^2 - 2r \cos\omega_0 z + r^2} = \frac{z^{-1}}{1 - 2r \cos\omega_0 z^{-1} + r^2 z^{-2}}, \quad |z| > r.
\end{equation}

\subsubsection*{c)}
\begin{itemize}
    \item \textbf{Comparison:} $X_1(z)$ and $X_2(z)$ share identical pole and zero locations in the complex $z$-plane, as well as an identical Region of Convergence ($|z| > r$). However, they are \textbf{not strictly identical} algebraically because they differ by an amplitude scale factor $G$:
    \begin{equation}
        X_1(z) = G \cdot X_2(z), \qquad \text{where } G = r \sin\omega_0.
    \end{equation}
    A pole-zero plot uniquely specifies a rational transfer function only up to a multiplicative constant.
    \item \textbf{Method to Derive $X_1(z)$ from the Pattern:} To resolve the unknown gain constant $K = G$ from the pole-zero geometry, one additional time-domain or frequency-domain constraint is required:
    \begin{enumerate}
        \item \textbf{Initial Value Property:} By the initial value theorem, the first sample is $x(0) = \lim_{z \to \infty} X(z) = 0$, which is consistent with the zero at $z=0$. Evaluating the second time sample via the limit theorem:
        \begin{equation}
            x(1) = \lim_{z \to \infty} z X(z) = \lim_{z \to \infty} \frac{K z^2}{z^2 - 2r\cos\omega_0 z + r^2} = K.
        \end{equation}
        From the known time-domain expression, $x_1(1) = r^1 \sin(\omega_0 \cdot 1) = r \sin\omega_0$. Setting the limit equal to $x_1(1)$ uniquely fixes the gain:
        \begin{equation}
            K = r \sin\omega_0 \implies X_1(z) = (r \sin\omega_0) X_2(z).
        \end{equation}
        \item \textbf{Frequency-Domain Evaluation:} Alternatively, measuring or specifying the system frequency response $X(e^{j\omega})$ at a known reference digital frequency $\omega = \omega_0$ allows direct algebraic solution for $K$.
    \end{enumerate}
\end{itemize}

\newpage
\section{Conclusion}
This laboratory study provided thorough theoretical and computational grounding in the $\mathcal{Z}$-transform, inverse $\mathcal{Z}$-transform, and discrete-time system analysis:
\begin{itemize}
    \item Linear convolution in the time domain was mapped to polynomial and rational multiplications in the transform domain, confirming that the output Region of Convergence contains the intersection of the individual input ROCs.
    \item First- and second-order difference equations were solved systematically for impulse, step, zero-input, and zero-state responses, demonstrating how poles govern exponential decay, oscillation frequency, and stability.
    \item Key algebraic properties (conjugation, real/imaginary parts, complex modulation, accumulation) were derived with rigorous region-of-convergence bounds.
    \item Ten distinct rational transforms from Proakis \& Manolakis Problem 3.14 were inverted into their causal time-domain sequences via partial fraction expansion and trigonometric decomposition.
    \item The parametric damped sinusoidal sequence in Problem 3.20 elucidated the fundamental duality between pole-zero geometry and time-domain dynamics, proving that pole-zero patterns define system modes while initial conditions or initial values calibrate the overall gain.
\end{itemize}

\end{document}
